\documentclass[11pt,oneside,openany]{book}
\usepackage[T1]{fontenc}
\usepackage{lmodern,microtype}
\usepackage[a4paper,margin=28mm]{geometry}
\usepackage{amsmath,amssymb,amsthm,mathtools}
\usepackage[hidelinks,bookmarksopen,bookmarksnumbered]{hyperref}
\hypersetup{pdftitle={Foundations of Numerical Analysis},
  pdfauthor={MathGenome},pdfsubject={Unified mathematical notes, P2-06}}
\allowdisplaybreaks[1]
\setlength{\emergencystretch}{2em}

\theoremstyle{plain}
\newtheorem{theorem}{Theorem}[chapter]
\newtheorem{lemma}[theorem]{Lemma}
\newtheorem{proposition}[theorem]{Proposition}
\newtheorem{corollary}[theorem]{Corollary}
\theoremstyle{definition}
\newtheorem{definition}[theorem]{Definition}
\newtheorem{example}[theorem]{Example}
\theoremstyle{remark}
\newtheorem{remark}[theorem]{Remark}

\begin{document}
\frontmatter
\hypersetup{pageanchor=false}
\begin{titlepage}
\thispagestyle{empty}
\vspace*{27mm}
{\large MathGenome \enspace / \enspace P2-06\par}
\vspace{8mm}
\hrule height 0.8pt
\vspace{11mm}
{\Huge\bfseries Foundations of\\[3pt] Numerical Analysis\par}
\vspace{4mm}
{\LARGE Error, Stability, and Certified Algorithms\par}
\vspace{12mm}
{\large Unified mathematical notes\par}
\vspace{6mm}
{\normalsize Floating-point error, conditioning, and stability\\
Matrix factorizations, least squares, and spectral perturbation\\
Interpolation, quadrature, root finding, and Newton methods\par}
\vfill
{\large 8 November 2026\par}
\end{titlepage}
\hypersetup{pageanchor=true}
\begingroup
\small
\setlength{\baselineskip}{10.6pt}
\tableofcontents
\endgroup

\chapter*{Scope and conventions}
\addcontentsline{toc}{chapter}{Scope and conventions}

The first three chapters develop the common language of rounding,
condition numbers, backward error, and finite-time stability.
Chapters 4 and 5 use this language for matrix factorizations,
least squares, eigenvalues, and low-rank approximation.
Chapters 6 and 7 treat interpolation, numerical differentiation,
and quadrature. Chapters 8 and 9 develop root enclosures, fixed
points, and local and semilocal Newton convergence.
The result map records the seventy-eight writing groups and
their relationships to the fifty original contracts.

A computed value in a rounding analysis is the exact real value
stored by the stated arithmetic model. Each primitive is subject to
its own safe-domain assumption. When an algorithm is studied
in exact arithmetic, its inputs and intermediates use the exact
operations of the indicated field. In a certification theorem,
a residual or an error enclosure is used only with the hypotheses
that make it a bound on the true quantity.

The established mathematical prerequisites include real calculus
from \cite{d03}, multivariable derivatives from \cite{d04},
finite-dimensional spectral theory and matrix factorizations
from \cite{d05}, the finite Chinese remainder theorem from
\cite{d07}, complete-space fixed points from \cite{d10}, and
the doubly stochastic decomposition from \cite{p205}.
Where a result is used as a dependency, the exact mathematical
interface is given in the text or companion map. The proof
of each new numerical bound gives the steps needed to transfer
the relevant exact theorem to the algorithm or certificate.
Textbook citations locate the statements and parallel treatments.

\mainmatter
\chapter{Floating-Point Models and Error Enclosures}\label{p206:ch01}

\section{The arithmetic model}

The rounding model used here concerns finitely many real numbers.
A stored number is identified with its real value. Exact arithmetic
inside an expression is indicated by ordinary arithmetic symbols, and
a rounding operation is indicated by \(\operatorname{fl}\).
The conventions follow \cite[Sections 2.1--2.5]{higham}.

\begin{definition}[A finite radix model]\label{p206:W01}
Fix an integer radix \(\beta\ge2\), an integer precision \(p\ge1\),
and integer exponent limits \(e_{\min}\le e_{\max}\).
The positive normal numbers are \(m\beta^e\), where
\(\beta^{p-1}\le m<\beta^p\) is an integer and
\(e_{\min}\le e\le e_{\max}\). Adjoin their negatives and zero.
With gradual underflow, also adjoin the signed numbers
\(m\beta^{e_{\min}}\) for \(1\le m<\beta^{p-1}\).
Let \(F\) be the resulting finite set.
Nearest rounding chooses a point of \(F\) at minimum distance,
with a fixed deterministic tie rule, for example the smaller
point in case of a tie. Set
\[
 u=\tfrac12\beta^{1-p},\qquad
 a_{\min}=\beta^{p-1+e_{\min}},\qquad
 a_{\max}=(\beta^p-1)\beta^{e_{\max}}.
\]
The normal safe domain used for relative rounding consists of zero
and the reals with \(a_{\min}\le |x|\le a_{\max}\).
\end{definition}

\begin{proposition}[Nearest rounding on the normal safe domain]
For every \(x\) in the normal safe domain,
\(\operatorname{fl}(x)=x(1+\delta)\) for some real \(\delta\)
with \(|\delta|\le u\).
\end{proposition}
\begin{proof}
Zero is represented exactly. For \(x>0\) in the safe range,
choose a binade
\([\beta^{p-1}\beta^e,\beta^p\beta^e]\) containing \(x\).
Its consecutive points are spaced by \(\beta^e\), including the
last gap up to the next binade when that endpoint exists.
In the top binade the assertion applies only up to \(a_{\max}\).
Thus one of the adjacent points is within \(\beta^e/2\) of \(x\),
and division by \(x\ge\beta^{p-1}\beta^e\) gives the bound.
For negative \(x\), the available distances are those of \(-x\)
under reflection of the representable set. The distance bound
is independent of the tie rule.
\end{proof}

The mathematical nearest map is defined on all reals. The normal
safe domain specifies where its relative-error theorem applies.
Other format models can use a smaller proved domain or a different
absolute bound near underflow.

\begin{definition}[Rounded operations]\label{p206:W02}
For representable inputs, rounded addition, subtraction,
multiplication and division are
\(\operatorname{fl}(a+b)\), \(\operatorname{fl}(a-b)\),
\(\operatorname{fl}(ab)\), and \(\operatorname{fl}(a/b)\),
respectively. Division requires \(b\ne0\).
A fused multiply-add is \(\operatorname{fl}(ab+c)\), with the
product and addition exact before the single rounding.
A rounded square root, when used, acts on a nonnegative
exact argument and applies one rounding to its real square root.
\end{definition}
For these models the relative estimate is a statement about
the exact result of each primitive. In a separate multiply and add,
both the product and the addition require their own range check.
For a fused expression, the range check is on the exact \(ab+c\).
The analysis below may equivalently start with an abstract relative
model on a specified domain if these same bounds are assumed.

\begin{proposition}[One-round error and enclosures]\label{p206:W03}
Suppose a computed \(z\) satisfies \(z=x(1+\delta)\),
\(|\delta|\le u<1\). Then \(x\) belongs to the closed interval
whose endpoints are \(z/(1-u)\) and \(z/(1+u)\), ordered by size.
In particular \(z=0\) implies \(x=0\) in this model. If instead
\(|z-x|\le\eta\), the interval \([z-\eta,z+\eta]\) contains \(x\).
If directed values \(z_\downarrow\le x\le z_\uparrow\)
are available, they themselves give an enclosure.
\end{proposition}
\begin{proof}
All possible divisors \(1+\delta\) are positive and the function
\(s\mapsto z/s\) is monotone or constant on that interval. This
gives the ordered relative enclosure for either sign of \(z\).
The absolute enclosure is the rearranged absolute inequality.
For a finite ordered format, a directed value can be obtained as
the largest representable point below, or the smallest above,
provided the respective defining set is nonempty. These choices
prove the directed inequalities separately.
\end{proof}

\begin{theorem}[Gradual underflow and Sterbenz's lemma]\label{p206:W04}
Assume gradual underflow. Nearest rounding of any \(|x|\le a_{\min}\)
has absolute error at most \(\beta^{e_{\min}}/2\).
If \(a,b\) are positive representable numbers with
\(a/2\le b\le2a\), then \(a-b\) is representable.
The latter assertion also applies to like-signed negative inputs
by reflection.
\end{theorem}
\begin{proof}
From zero to \(a_{\min}\) the grid, including both endpoints,
has spacing \(\beta^{e_{\min}}\). Reflection handles the
negative side and nearest rounding gives the half-spacing bound.

For exact subtraction interchange the inputs and negate the
difference if necessary, so \(b\le a\le2b\).
Use the grid exponent \(e\) of \(b\), with \(e=e_{\min}\) if it is
subnormal. Write \(b=m\beta^e\), \(m\le\beta^p-1\).
Since \(a\ge b\), its grid exponent is at least \(e\).
Thus \(a-b=k\beta^e\) for an integer \(0\le k\le m\).
If \(k=0\) the result is zero. If \(0<k<\beta^{p-1}\),
successively multiply \(k\) by \(\beta\) and reduce \(e\) by one,
until the coefficient is normal or the minimum exponent is reached.
At the first normal step the coefficient is less than \(\beta^p\).
At the minimum exponent a shorter coefficient is allowed.
This proves representability.
\end{proof}
Exact subtraction of stored inputs is compatible with sensitivity
to their original errors. For \(a-b\ne0\), individual input
perturbations bounded by \(\epsilon |a|\) and \(\epsilon |b|\)
can change the difference by \(\epsilon(|a|+|b|)\).
Opposite signs for the two perturbations attain this bound.
The corresponding relative amplification may be large even
when the subtraction primitive is exact.

\begin{proposition}[Outward interval operations]\label{p206:W05}
Let \(X=[a,b]\) and \(Y=[c,d]\). Addition is enclosed by outward
rounding of \([a+c,b+d]\), subtraction by that of \([a-d,b-c]\),
and multiplication by that of the interval between the minimum
and maximum of \(ac,ad,bc,bd\). If \(0\notin Y\), the reciprocal
is enclosed by outward rounding of \([1/d,1/c]\), and division
is enclosed by multiplying with that reciprocal enclosure.
Each outward endpoint used must exist.
\end{proposition}
\begin{proof}
The affine dependence on the coordinates gives the two signed
addition bounds. For a product each coordinate separately has
an affine dependence, so its extrema over the rectangle occur
at corners. The reciprocal is decreasing on each component of
\(\mathbb R\setminus\{0\}\). These real ranges contain all exact
values, and replacing their left and right endpoints by a lower
and an upper bound preserves containment. The stated product
argument applies to the enlarged reciprocal interval.
\end{proof}

\section{Accumulation of rounding factors}

\begin{lemma}[Signed products of rounding factors]\label{p206:W06}
Let \(|\delta_i|\le u\) for \(1\le i\le n\), with \(u\ge0\)
and \(u<1\). If \(nu<1\) and \(\rho_i\in\{-1,1\}\), then
\[
 \prod_{i=1}^n(1+\delta_i)^{\rho_i}=1+\theta,\qquad
 |\theta|\le \gamma_n:=\frac{nu}{1-nu}.
\]
The empty product has \(\theta=0\).
\end{lemma}
\begin{proof}
For either sign of \(\rho_i\), the factor lies in
\([1-u,(1-u)^{-1}]\), using \(1+u\le(1-u)^{-1}\).
Induction gives \((1-u)^n\ge1-nu\), so the product is between
\(1-nu\) and \((1-nu)^{-1}\). The farther endpoint from one
is at distance \(\gamma_n\).
\end{proof}

\begin{theorem}[Sequential dot product]\label{p206:W07}
Suppose the inputs \(x_i,y_i\) are exact stored numbers and all
operations below satisfy the relative model with \(nu<1\).
Either form of the following sequential dot product produces
\[
 \widehat s=\sum_{i=1}^n x_i y_i(1+\theta_i),\qquad
 |\theta_i|\le\gamma_n,\qquad
 \left|\widehat s-\sum_{i=1}^n x_i y_i\right|
       \le\gamma_n\sum_{i=1}^n |x_i y_i|.
\]
The first form separately rounds each product and adds the stored
products from left to right, starting at the first one. The second
starts from zero and applies one fused multiply-add per term.
The empty dot product is zero.
\end{theorem}
\begin{proof}
In the first form, write the product rounding factor of term
\(i\) as \(1+\alpha_i\) and the sum rounding factor when term \(j\)
is inserted as \(1+\delta_j\), \(j\ge2\).
The first term has the product factor and \(n-1\) sum factors.
Term \(i\ge2\) has its product factor and sum factors for
\(j=i,\ldots,n\), at most \(n\) in total. This follows by
distributing each sum factor over the two terms at its step.
In the fused recurrence a term \(i\) acquires precisely the factors
for steps \(i,\ldots,n\). Lemma~\ref{p206:W06} bounds each
path, and the triangle inequality gives the absolute estimate.
\end{proof}
In a different evaluation order the bound must use a proved path
length \(N\) and its condition \(Nu<1\). In particular, a bound
using \(N=2n\) requires \(2nu<1\), even if a sharper analysis of
the same algorithm only needs \(nu<1\).

\begin{theorem}[Summation depth and matrix products]\label{p206:W08}
For a binary rounded sum of real inputs \(x_i\), let \(d\)
bound the number of additions from any input to the output.
Under the relative model and \(du<1\),
the absolute error is at most \(\gamma_d\sum_i|x_i|\).
For \(n\ge1\) leaves, a balanced tree can be chosen with
\(d=\lceil\log_2 n\rceil\). If the exact sum \(s\) is nonzero,
its componentwise relative condition number is
\(\sum_i|x_i|/|s|\).
For a matrix product with inner dimension \(k\), computing each
entry as in Theorem~\ref{p206:W07} under \(ku<1\) yields
\[
       |\widehat{AB}-AB|\le \gamma_k\,|A|\,|B|.
\]
Here matrix absolute values and inequalities are entrywise.
\end{theorem}
\begin{proof}
Expand the tree recursively. Each leaf has one factor per
rounded addition on its path, so Lemma~\ref{p206:W06} and
the triangle inequality give the estimate. Splitting any group
of leaves into two groups whose sizes differ by at most one gives
the asserted depth by induction. An additional independently
rounded product at each leaf increases the bound on depth by one.

If \(|\Delta x_i|\le\epsilon |x_i|\), the change of the exact sum
has modulus at most \(\epsilon\sum_i |x_i|\). Choosing
\(\Delta x_i=\epsilon|x_i|\) for every \(i\) attains the bound.
Divide by \(\epsilon|s|\). At \(s=0\) this argument yields an
absolute bound without dividing. Apply the dot-product estimate
to each entry to obtain the last formula, including \(k=0\)
with zero dot products.
\end{proof}
\chapter{Conditioning and Backward Error}

\section{Local sensitivity}

A condition number is assigned to a mathematical problem and a
choice of input and output scales. Throughout this chapter a matrix
norm used with a vector norm is subordinate to it. Different input
and output spaces may carry different vector norms.

\begin{definition}[Absolute and relative condition numbers]\label{p206:W09}
Let \(f:U\to Y\) be Fr\'echet differentiable at \(x\) in an open
subset \(U\) of a finite-dimensional normed space \(X\), with \(Y\)
also normed. Set
\[
 \kappa_{\mathrm{abs}}(f,x)=\|Df(x)\|,\qquad
 \kappa_{\mathrm{rel}}(f,x)
   =\|Df(x)\|\frac{\|x\|}{\|f(x)\|}\quad (x\ne0,\ f(x)\ne0).
\]
The latter definition uses the indicated relative scales.
At a zero scale the absolute condition number remains defined.
A singular problem that fails to define a differentiable solution
map lies outside this derivative formula.
\end{definition}
The first number has a worst-case interpretation. For sufficiently
small \(r>0\) define
\[
 K(r)=r^{-1}\sup_{\|v\|\le r}\|f(x+v)-f(x)\|.
\]
Differentiability gives a uniform remainder bound
\(\|f(x+v)-f(x)-Df(x)v\|\le\omega(r)\|v\|\) for
\(\|v\|\le r\), with \(\omega(r)\to0\).
Hence \(K(r)\le\|Df(x)\|+\omega(r)\). For any unit vector \(w\),
test \(v=rw\) to obtain \(K(r)\ge\|Df(x)w\|-\omega(r)\).
Taking the supremum in \(w\) gives convergence to
\(\|Df(x)\|\). If \(X\) is zero dimensional then \(K(r)=0\).
The relative interpretation follows by replacing the input radius
by \(r\|x\|\) and dividing the output by \(\|f(x)\|\).
See \cite[Sections 1.5--1.6]{higham} for these scales.

\begin{theorem}[Neumann series and inverse perturbation]\label{p206:W10}
Let \(A:X\to Y\) be a bounded invertible linear map of Banach
spaces with bounded inverse and let \(\Delta A:X\to Y\) be
bounded linear. If \(E=A^{-1}\Delta A\) satisfies
\(\|E\|<1\), then \(A+\Delta A\) is invertible and
\[
 \|(A+\Delta A)^{-1}\|\le \frac{\|A^{-1}\|}{1-\|E\|}.
\]
The inverse of \(I+E\) is \(\sum_{j\ge0}(-E)^j\), and
\[
 (A+\Delta A)^{-1}-A^{-1}
       =-(A+\Delta A)^{-1}\Delta A\,A^{-1}.
\]
\end{theorem}
\begin{proof}
The algebra of bounded operators on \(X\) is complete in its
operator norm. The series converges since
\(\|(-E)^j\|\le\|E\|^j\), including \(j=0\) with
\(\|I\|\le1\). Multiplying the partial sums on either side
by \(I+E\) gives \(I-(-E)^{N+1}\), which tends to \(I\).
Thus the sum is the inverse and has norm at most
\((1-\|E\|)^{-1}\). Factor \(A+\Delta A=A(I+E)\)
and apply the norm bound. For the difference identity, insert
the factors \(A\) and \(A+\Delta A\) between the two inverses and
subtract. The argument also works in a unital Banach algebra with
\(\|1\|\le1\). The geometric-series identity is available in
the fixed Mathlib as \texttt{Units.oneSub} with the sign changed
\cite{mathlib}.
\end{proof}

\begin{theorem}[Perturbed linear systems]\label{p206:W11}
Let \(A\) be invertible, \(Ax=b\), and
\(\|A^{-1}\|\|\Delta A\|<1\). Put
\((A+\Delta A)\tilde x=b+\Delta b\).
If \(b\ne0\), define
\(a=\|\Delta A\|/\|A\|\), \(c=\|\Delta b\|/\|b\|\),
and \(\kappa(A)=\|A\|\|A^{-1}\|\). Then
\[
 \frac{\|\tilde x-x\|}{\|x\|}
 \le \frac{\kappa(A)(a+c)}{1-\kappa(A)a}.
\]
For any \(b\), the corresponding absolute bound is
\[
 \|\tilde x-x\|
 \le \frac{\|A^{-1}\|}{1-\|A^{-1}\|\|\Delta A\|}
       \bigl(\|\Delta b\|+\|\Delta A\|\|x\|\bigr).
\]
\end{theorem}
\begin{proof}
The preceding theorem gives the perturbed inverse. Subtract
\((A+\Delta A)x=b+\Delta A x\) from the perturbed equation to get
\((A+\Delta A)(\tilde x-x)=\Delta b-\Delta A x\).
Taking norms gives the absolute estimate.
If \(b\ne0\), then \(x\ne0\) and
\(\|\Delta b\|/\|x\|\le c\|A\|\). Divide and simplify.
The case \(b=0\) uses the absolute version.
\end{proof}

\section{Forward and backward errors}

\begin{definition}[Two error measurements]\label{p206:W12}
For exact data \(x\) and an observed output \(\hat y\), the
absolute forward error is \(\|\hat y-f(x)\|\). Its relative
version divides by the chosen nonzero output scale.
For backward error choose an admissible set of perturbations
\(\mathcal D_x\) and a nonnegative size \(s(v)\) on that set.
The backward error is
\[
 \inf\{s(v):v\in\mathcal D_x,\ x+v\in U,\ f(x+v)=\hat y\},
\]
with the infimum \(+\infty\) when there is no feasible
perturbation. A backward certificate is a particular feasible
perturbation together with a bound on its size.
\end{definition}
For example one can use the norm of a perturbation, its
relative norm, or componentwise bounds scaled by selected data.
A finite infimum guarantees the existence of arbitrarily close
feasible sizes. To guarantee a minimum one can assume in finite
dimension that the feasible perturbations form a nonempty closed set
and that the size is a norm. A minimizing sequence then has a
bounded subsequence converging to a feasible minimizer.
The closedness assumption matters, for example when admissible
inputs are required to stay in an open region.

\begin{theorem}[Conditioning converts backward error to forward error]
\label{p206:W13}
Suppose \(f\) is Fr\'echet differentiable at \(x\), with
\(x,f(x)\ne0\). For actual perturbations \(v\) with
\(\hat y=f(x+v)\) and \(\epsilon=\|v\|/\|x\|\to0\),
\[
 \frac{\|\hat y-f(x)\|}{\|f(x)\|}
    \le \kappa_{\mathrm{rel}}(f,x)\epsilon+o(\epsilon).
\]
For a fixed segment \([x,x+v]\), continuity on the segment,
differentiability along its interior, and a uniform derivative
bound \(\|Df\|\le L\) there give
\(\|f(x+v)-f(x)\|\le L\|v\|\).
\end{theorem}
\begin{proof}
Write the Fr\'echet remainder and divide its estimate by
\(\|f(x)\|\). Its modulus of smallness can be taken uniform
for all perturbations in a sufficiently small ball, as in
the proof following Definition~\ref{p206:W09}.
If the reported relative backward error is an infimum \(b_*\)
rather than an attained perturbation, use perturbations of size
at most \(b_*+\tau\). For each fixed \(\epsilon_0>0\), the
ratio is bounded by
\((\kappa_{\mathrm{rel}}+\epsilon_0)(b_*+\tau)\)
whenever these perturbations are sufficiently small. Letting
\(\tau\downarrow0\) gives the same first-order conclusion for
the infimum, including zero by continuity.

For the finite bound apply the real mean value theorem to
a norm-one real functional of \(t\mapsto f(x+tv)\) that
attains the norm on the difference of its endpoint values.
Such a functional is given by Hahn--Banach, including on the
realification of a complex space. Its scalar derivative is
bounded in modulus by \(L\|v\|\), which proves the assertion.
If the endpoint difference is zero the bound holds directly.
This segment form is also the estimate in
\texttt{Convex.norm\_image\_sub\_le\_of\_norm\_fderiv\_le}
of the fixed Mathlib \cite{mathlib}; the Euclidean version
is in the multivariable notes \cite{d04}.
\end{proof}

\begin{theorem}[Linear residual certificates]\label{p206:W14}
Let \(A\) be invertible, \(Ax=b\), and \(r=b-A\hat x\).
Then
\[
 \|\hat x-x\|\le\|A^{-1}\|\|r\|,\qquad
 \frac{\|\hat x-x\|}{\|x\|}
      \le \kappa(A)\frac{\|r\|}{\|b\|}\quad (b\ne0).
\]
The absolute normwise backward error when only the right-hand side
can change is exactly \(\|r\|\).
\end{theorem}
\begin{proof}
The identity \(x-\hat x=A^{-1}r\) gives the first bound,
and \(\|b\|\le\|A\|\|x\|\) gives the second. The equation
\(A\hat x=b+\Delta b\) forces \(\Delta b=-r\).
Thus the backward certificate exists, is unique in this model,
and has the claimed size.
\end{proof}

\begin{theorem}[Joint normwise and componentwise backward error]
\label{p206:W15}
Given \(A,b,\hat x\), let \(r=b-A\hat x\).
For \(\alpha,\beta\ge0\), allow
\(\|\Delta A\|\le\epsilon\alpha\) and
\(\|\Delta b\|\le\epsilon\beta\), where the matrix norm is
subordinate. The least \(\epsilon\) for which
\((A+\Delta A)\hat x=b+\Delta b\) is possible is
\[
                 \frac{\|r\|}{\alpha\|\hat x\|+\beta}.
\]
For nonnegative scales \(E\) and \(f\) on entries of \(A\)
and \(b\), the analogous componentwise answer is
\[
     \max_i\frac{|r_i|}{(E|\hat x|+f)_i}.
\]
A zero numerator and denominator give zero, a nonzero numerator
over zero gives \(+\infty\), and the maximum over no rows is zero.
Finite values are attained in each model.
\end{theorem}
\begin{proof}
The equation \(r=\Delta A\hat x-\Delta b\) and the triangle
inequality give both lower bounds. For the normwise bound let
\(D=\alpha\|\hat x\|+\beta>0\) and put \(c=r/D\).
When \(\hat x\ne0\) choose a dual functional \(\phi\) of norm one
with \(\phi(\hat x)=\|\hat x\|\) by Hahn--Banach.
Take \(\Delta A=\alpha c\otimes\phi\), acting as
\(v\mapsto\alpha c\,\phi(v)\), and take \(\Delta b=-\beta c\).
If \(\hat x=0\), take \(\Delta A=0\) and the same \(\Delta b\).
When \(D=0\) only \(r=0\) is feasible, attained by zero perturbations.

For the componentwise bound put \(D_i=(E|\hat x|+f)_i\).
When \(D_i>0\) set \(s_i=r_i/D_i\).
Choose scalars \(\zeta_j\) with \(|\zeta_j|\le1\) and
\(\zeta_j\hat x_j=|\hat x_j|\), for example the conjugate
phase when \(\hat x_j\ne0\) and zero otherwise.
Use \(\Delta A_{ij}=s_i E_{ij}\zeta_j\) and
\(\Delta b_i=-s_i f_i\). This produces exactly \(r_i\)
with size at most the proposed maximum. In a row with \(D_i=0\),
a nonzero residual is infeasible; for a zero residual set
the perturbation row to zero. This proves attainment for all
finite cases. See \cite[Sections 7.1--7.2]{higham}.
\end{proof}
\chapter{Stability, Consistency, and Convergence}

\section{Finite-horizon definitions}

Let \(X\) be a Banach space. For every \(h>0\) let \(C_h\) be
a bounded linear operator on \(X\). Let \(T(t)\), \(t\ge0\), be a
strongly continuous semigroup of bounded linear operators, so
\(T(0)=I\), \(T(s+t)=T(s)T(t)\), and \(t\mapsto T(t)x\) is
continuous for every \(x\in X\). Thus \(C_h^n\) and \(T(nh)\)
act on the same space. Write \(T_0>0\) for a fixed horizon.
The distinctions below specialize the consistency and stability
framework in \cite[Section 4.2]{trefpde}.

\begin{definition}[Stability, defects and convergence]\label{p206:W16}
Finite-horizon small-step stability means that for this \(T_0\)
there are \(h_0>0\) and \(M<\infty\) with
\[
 \|C_h^n\|\le M\quad\text{for all }0<h\le h_0,\quad
                        n\ge0,\quad nh\le T_0.
\]
Pointwise local consistency at \(x\) means
\(h^{-1}\|(C_h-T(h))x\|\to0\).
Uniform orbit consistency for \(x\) on this horizon means
\[
 \tau_h(x):=\max_{\substack{j\ge0\\(j+1)h\le T_0}}
      h^{-1}\|(C_h-T(h))T(jh)x\|\longrightarrow0,
\]
where a maximum over no steps is zero.
Strong convergence, uniform in mesh times on this horizon, means
\[
 E_h(x):=\max_{\substack{n\ge0\\nh\le T_0}}
                    \|C_h^n x-T(nh)x\|\longrightarrow0
              \quad\text{for each }x\in X.
\]
For statements on all finite horizons the bounds may depend on \(T_0\).
\end{definition}
For \(x\) in the generator domain, the one-step pointwise
condition is equivalent to
\((C_h-I)x/h\to Ax\). Indeed, by definition of the generator,
\((T(h)-I)x/h\to Ax\), and their difference is the local defect
with its sign. This equivalence is pointwise.
The positive convergence theorem below uses the stronger control
along the trajectory, or assumes that this control has been
verified on a dense class.

\section{Error propagation and stability}

\begin{theorem}[A finite-time error recurrence]\label{p206:W17}
Let \(L,\tau_h\ge0\), \(h>0\), and let nonnegative errors satisfy
\[
 e_{j+1}\le(1+Lh)e_j+h\tau_h
                 \quad\text{whenever }(j+1)h\le T_0.
\]
For every \(nh\le T_0\),
\[
 e_n\le e^{LT_0}e_0+T_0e^{LT_0}\tau_h .
\]
If \(e_0,\tau_h\to0\) as \(h\downarrow0\), the maximum of \(e_n\)
on the horizon tends to zero. More generally, a nonnegative
\(\tau_{h,j}\) in the recurrence contributes
\(\sum_{j=0}^{n-1}h(1+Lh)^{n-1-j}\tau_{h,j}\).
\end{theorem}
\begin{proof}
Successive substitution, or induction on \(n\), gives the weighted
sum and the initial term \((1+Lh)^n e_0\).
Because \(1+Lh\le e^{Lh}\) and \(nh\le T_0\), every weight
without its factor \(h\) is at most \(e^{LT_0}\).
The sum of the \(n\) factors \(h\) is at most \(T_0\).
This proof includes \(L=0\) and \(n=0\).
\end{proof}

We use the following form of uniform boundedness. If a family of
bounded linear maps on a Banach space is bounded at each \(x\),
then its operator norms are uniformly bounded. For completeness,
the sets \(F_m=\{x:\sup_i\|S_i x\|\le m\}\) are closed and cover
the space. By the Baire category theorem one contains an open
ball. If a ball of radius \(r\) about \(x_0\) lies in \(F_m\),
then \(\|S_i v\|\le2m\) for \(\|v\|\le r/2\) by subtracting
the values at \(x_0+v\) and \(x_0\). Scaling yields the common
bound \(4m/r\). Baire and the uniform-boundedness theorem are
available in the fixed Mathlib \cite{mathlib}.
In particular, strong continuity on \([0,T_0]\) bounds each
orbit, so
\[
                   K_{T_0}:=\sup_{0\le t\le T_0}\|T(t)\|<\infty.
\]

\begin{theorem}[Convergence implies small-step stability]
\label{p206:W18}
Strong convergence uniform in mesh times on \([0,T_0]\) for
every initial vector implies finite-horizon small-step stability
on \([0,T_0]\).
\end{theorem}
\begin{proof}
Suppose there is no common operator bound on any sufficiently
small step range. For every \(j\ge1\) choose \(0<h_j\le 1/j\)
and \(n_jh_j\le T_0\) with \(\|C_{h_j}^{n_j}\|>j\).
For every \(x\in X\), the difference between \(C_{h_j}^{n_j}x\)
and \(T(n_jh_j)x\) tends to zero. Since the exact orbit is bounded,
the family \(C_{h_j}^{n_j}\) is bounded at each \(x\).
Uniform boundedness gives a common operator norm bound, a
contradiction.
\end{proof}

\begin{example}[The step restriction]
On \(X=\mathbb R\) take \(T(t)=I\) and \(T_0=1\).
For \(0<h\le1/2\), let \(C_h=I\). For \(h>1/2\) with
\(h\ne3/4\), let \(C_h\) be scalar multiplication by
\(1/|h-3/4|\). Define \(C_{3/4}=I\).
All operators are individually bounded and small-step convergence
is exact, while \(n=1\) and \(h\to3/4\) prevent a common bound
over every \(h\) with \(nh\le1\).
\end{example}

\begin{theorem}[Stability and uniform orbit consistency imply convergence]
\label{p206:W19}
Under finite-horizon small-step stability, for every \(x\)
and \(0<h\le h_0\),
\[
               E_h(x)\le MT_0\tau_h(x).
\]
If \(\tau_h(x)\to0\) on a dense class of \(x\), then
\(E_h(x)\to0\) for all \(x\in X\).
\end{theorem}
\begin{proof}
For \(n\ge1\) the identity
\[
 C_h^nx-T(nh)x
     =\sum_{j=0}^{n-1}C_h^{n-1-j}(C_h-T(h))T(jh)x
\]
follows by telescoping consecutive terms
\(C_h^{n-j}T(jh)x\). It uses the semigroup property and requires
no commutation between \(C_h\) and \(T(h)\).
When \(nh\le T_0\), each norm is bounded by \(Mh\tau_h(x)\).
Summing and including \(n=0\) proves the estimate.
For an arbitrary \(x\), approximate by \(y\) from the dense
class. For each step and time in the same range,
\[
  E_h(x)\le E_h(y)+(M+K_{T_0})\|x-y\|.
\]
First let \(h\downarrow0\) for the fixed \(y\), then
let \(\|x-y\|\) tend to zero.
\end{proof}
\chapter{Factorizations and Least Squares}
\label{p206:factorizations}

A superscript \(^{*}\) denotes conjugate transpose, or transpose over
\(\mathbb R\). Absolute values of matrices in entrywise inequalities
are taken entry by entry, and their products are ordinary matrix
products of nonnegative entries. Matrix norms with subscript \(2\)
are operator norms for the Euclidean norm. The exact arguments
allow real or complex data where specified. The rounding bounds
use the real scalar model of Chapter~1, with no exceptional range
events in any listed operation. Input entries are the given stored
numbers. The notation \(\gamma_k=ku/(1-ku)\) and the signed product
bound are as in Lemma~\ref{p206:W06}.

\section{Elimination and triangular solves}

\begin{theorem}[Exact row-pivoted LU]\label{p206:W20}
Every square matrix over a field admits \(PA=LU\), where \(P\) is a
permutation matrix, \(L\) is unit lower triangular, and \(U\) is in
row-echelon form and is upper triangular. Elimination may skip
a column whose remaining entries are all zero.
For elimination without row exchanges and with a pivot in every
column up to the last, nonzero division pivots are equivalent to
nonvanishing leading principal minors of orders \(1,\ldots,n-1\).
The last diagonal entry is nonzero exactly when the full determinant
is nonzero, provided the earlier pivots exist.
\end{theorem}

\begin{proof}
After \(r\) pivots, keep \(PA=LB\), where \(B\) is the current matrix
and \(L\) is unit lower triangular with off-diagonal entries supported
in its first \(r\) columns. Initially \(P=L=I\), \(B=A\). Scan the
columns from left to right. If the current column is zero on rows
\(r+1,\ldots,n\), advance the column and keep the row cursor \(r+1\).
Otherwise choose a nonzero entry among those rows and exchange its
row with row \(r+1\). If the exchange is \(S\), replace \(P,B,L\)
by \(SP,SB,SLS^{-1}\), respectively. The new \(L\) is still of the
prescribed form since this exchange fixes its first \(r\) indices.
Concretely, the exchange must include the entries of the previous
multiplier columns in those rows.

Divide the remaining entries in the pivot column by the pivot and
store the quotients as the tail of a vector \(\ell\) supported on
rows \(r+2,\ldots,n\). Subtract each multiplier times the pivot row
from its row. This multiplies \(B\) by \(I-\ell e_{r+1}^{\mathsf T}\),
where \(e_{r+1}\) is a coordinate vector. Its inverse is
\(I+\ell e_{r+1}^{\mathsf T}\) because \(e_{r+1}^{\mathsf T}\ell=0\).
Multiplying \(L\) by this inverse stores \(\ell\) in its next multiplier column,
since \(L\ell=\ell\). These operations preserve every previously
cleared or skipped column on the unprocessed rows. At the end the
pivot columns increase strictly with the pivot row, and all rows
after the last pivot are zero. The \(i\)-th pivot column is at
least \(i\), which gives the upper-triangular form. The empty matrix
has the empty factorization.

For the unexchanged algorithm, the elimination steps through
pivot \(k-1\) reduce the leading \(k\)-square block by
determinant-preserving row additions. Its determinant is
the product of the first \(k\) pivots as soon as the preceding
\(k-1\) pivots have been computed. Induction gives the claimed
criterion and the assertion about the terminal diagonal. An
invertible matrix can have a zero leading minor, as shown by
the matrix exchanging two coordinates. Existence of a bare
unpivoted factorization is a weaker assertion. For example
\(\operatorname{diag}(0,1)=I\operatorname{diag}(0,1)\), although a
Doolittle step dividing by its first diagonal entry is unavailable.
\end{proof}

For a fixed sequential subtraction kernel let
\[
  s_t=\operatorname{fl}\bigl(s_{t-1}
                    -\operatorname{fl}(a_tb_t)\bigr),
  \qquad 1\le t\le k .
\]
If the multiplication and subtraction at step \(t\) have factors
\(1+\mu_t\) and \(1+\sigma_t\), the real identities are
\[
 \begin{aligned}
 s_{t-1}&=\frac{s_t}{1+\sigma_t}+a_tb_t(1+\mu_t),\\
 s_0&=s_k\prod_{j=1}^{k}(1+\sigma_j)^{-1}
      +\sum_{t=1}^k a_tb_t(1+\mu_t)
                      \prod_{j<t}(1+\sigma_j)^{-1}.
 \end{aligned}                                                   \tag{4.1}\label{p206:subtracttrace}
\]
The second identity follows by successively substituting the first.
An individual summand uses \(t\) factors, regardless of the total
number of operations. When a final division or square root is
undone, its factors affect only the terminal term of
\eqref{p206:subtracttrace}.

\begin{theorem}[Componentwise backward error of a triangular solve]
\label{p206:W21}
Let \(T\in\mathbb R^{n\times n}\) be triangular with nonzero diagonal.
Use forward or backward substitution with sequential rounded
multiply-subtract steps and one final rounded division in each row.
Suppose every operation is safe and \(nu<1\). The result
\(\hat x\) satisfies
\[
        (T+\Delta T)\hat x=b,\qquad
        |\Delta T|\le \gamma_n|T|.
\]
The correction preserves the triangular pattern and all structural
zeros.
\end{theorem}

\begin{proof}
In row \(i\) begin with \(s_0=b_i\) and apply the sequential
kernel to pairs \(T_{ij},\hat x_j\) for already determined
coordinates \(j\). If there are \(k\) such entries,
\[
 s_k=T_{ii}\hat x_i(1+\delta_i)^{-1}.
\]
By \eqref{p206:subtracttrace} each off-diagonal coefficient \(T_{ij}\) is multiplied
by at most \(k\) signed factors and the diagonal coefficient
by \(k+1\le n\). Apply Lemma~\ref{p206:W06} to each coefficient
and collect the perturbed coefficients row by row. A zero
coefficient stays zero under this description, including the
case of a zero product. When \(n=0\) the assertion is empty.
\end{proof}

\begin{theorem}[A safe sequential Doolittle run]\label{p206:W22}
Consider real square elimination, with row exchanges as in
Theorem~\ref{p206:W20}. Assume columns \(1,\ldots,n-1\) are
processed in order without skips, all pivots used as divisors
are nonzero, and each division and trailing update uses the
relative error model.
Store each multiplier by a rounded division and update each
trailing entry by the sequential multiply-subtract step
\eqref{p206:subtracttrace}.
If \(3nu<1\), the stored factors satisfy
\[
 PA+\Delta A=\hat L\hat U,\qquad
             |\Delta A|\le\gamma_{3n}|\hat L|\,|\hat U|.
\]
\end{theorem}

\begin{proof}
Use final row positions to label the stored quantities.
A future row exchange relabels its trailing data and its already
computed multipliers together. Therefore the row eventually
placed at position \(i\) is updated with the multiplier
\(\hat l_{it}\) at each preceding pivot \(t<i\).
Fix an entry \((i,j)\) and put \(k=\min(i,j)\). Its successive
values from \(s_0=(PA)_{ij}\) use products \(\hat l_{it}\hat u_{tj}\)
for \(t=1,\ldots,k-1\). If \(i\le j\), the value after these updates
is stored as \(\hat u_{ij}\), with \(\hat l_{ii}=1\).
If \(i>j\), the multiplier division gives
\[
       s_{j-1}=\hat l_{ij}\hat u_{jj}(1+\delta_{ij})^{-1}.
\]
Hence \eqref{p206:subtracttrace} expresses \((PA)_{ij}\) as the sum of its exact
\((\hat L\hat U)_{ij}\) terms, each multiplied by a product of
at most \(n\) signed factors. The terminal product uses at most
\(k\) factors and earlier terms use at most \(k-1\).
The difference from the unweighted sum is bounded by
\(\gamma_n\sum_t|\hat l_{it}\hat u_{tj}|\).
Since \(\gamma_n\le\gamma_{3n}\), this yields the advertised
conservative constant. The empty factorization again satisfies
the statement. The trace bound is conditional on the specified
run, including its pivot hypotheses.
\end{proof}

\begin{theorem}[Partial pivot growth and an end-to-end solve certificate]
\label{p206:W23}
For exact elimination on a nonzero \(n\)-square real matrix,
choose a largest-magnitude available entry in the current pivot
column. The maximum magnitude among all entries in the successive
trailing submatrices and computed \(U\) rows is at most
\(2^{n-1}\max_{ij}|a_{ij}|\). Choosing the first row on a tie can
attain the bound.

For a rounded run satisfying Theorem~\ref{p206:W22}, suppose
the forward and backward triangular solves of
\(\hat L y=Pb\), \(\hat Ux=y\) satisfy Theorem~\ref{p206:W21}.
Then the computed \(\hat x\) solves
\[
 (PA+\Delta A_{\rm solve})\hat x=Pb,\qquad
 |\Delta A_{\rm solve}|
   \le(\gamma_{3n}+2\gamma_n+\gamma_n^2)|\hat L|\,|\hat U|.
\]
\end{theorem}

\begin{proof}
With exact partial pivoting \(|l_{it}|\le1\). Every Schur update
has the form \(a_{ij}^{\rm new}=a_{ij}^{\rm old}-l_{it}u_{tj}\),
so the old maximum on the active matrix bounds its magnitude
by twice that maximum. Row exchanges preserve maxima. At most
\(n-1\) such growth steps occur. For sharpness use the matrix whose
columns \(j<n\) have \(1\) on the diagonal, \(-1\) below it, and
\(0\) above it, with all entries in the last column equal to \(1\).
No pivot exchange is made under the stated tie rule. All
multipliers below each of the first \(n-1\) pivots are \(-1\).
After each update the remaining last-column entries double,
while the later columns other than the last retain their
initial form. The last diagonal entry is \(2^{n-1}\).
The growth bound concerns exact partial pivoting. The solve
certificate uses the actual stored factors.

By the two triangular certificates,
\[
 (\hat L+\Delta L)\hat y=Pb,\qquad
 (\hat U+\Delta U)\hat x=\hat y .
\]
Expand the product of the two perturbed factors. Its correction
to \(PA\) is
\(\Delta A+\Delta L\hat U+\hat L\Delta U+\Delta L\Delta U\).
Entrywise triangle inequalities and
\(|\Delta L|\le\gamma_n|\hat L|\),
\(|\Delta U|\le\gamma_n|\hat U|\) give the bound. To obtain an
error in the original system, multiply by \(P^{-1}\), and use
the inverse or conditioning bounds of Chapter~2 when the
resulting normwise perturbation satisfies their hypotheses.
\end{proof}

These elimination and triangular solve conventions are comparable
to \cite[Chapters~8--9]{higham} and \cite[Sections~3.2--3.4]{golub}.

\section{Positive matrices and Cholesky factorization}

\begin{theorem}[Positive-definite Cholesky]\label{p206:W24}
A real symmetric or complex Hermitian positive-definite matrix
\(A\) has a unique factorization \(A=R^*R\) with \(R\) upper
triangular and positive real diagonal. Equivalently
\(A=LL^*\), where \(L=R^*\) is lower triangular with positive real
diagonal.
\end{theorem}

\begin{proof}
Use induction, starting with the empty matrix. Partition
\[
 A=\begin{pmatrix}\alpha&a^*\\a&B\end{pmatrix}.
\]
The first diagonal entry \(\alpha\) is positive. In any factor
of the stated form the first row must be
\((\sqrt{\alpha},a^*/\sqrt{\alpha})\), by comparing the first
row of \(R^*R\) with that of \(A\). The remaining factor must
factor \(B-aa^*/\alpha\). This Schur complement is positive
definite because for every nonzero \(y\) its quadratic form
equals that of \(A\) at \((-a^*y/\alpha,y)\), which is nonzero.
Induction gives the unique remaining block, and block
multiplication proves the factorization.
\end{proof}

\begin{theorem}[A canonical semidefinite Cholesky factor]
\label{p206:W25}
A Hermitian positive-semidefinite matrix \(A\) has a unique factor
\(A=R^*R\) among upper triangular \(R\) whose diagonal is real
nonnegative and whose entire \(i\)-th row vanishes whenever
\(R_{ii}=0\).
\end{theorem}

\begin{proof}
If the leading diagonal is positive, repeat the block argument
of Theorem~\ref{p206:W24}. The Schur complement is now
positive semidefinite by the same quadratic-form identity. If
the leading diagonal is zero, then each off-diagonal entry in
its row is zero. Indeed, for the block above, if \(a^*y\ne0\),
the quadratic form at \((t,y)\) equals
\(2\operatorname{Re}(\overline{t}\,a^*y)+y^*By\) and becomes
negative on choosing \(t=-s a^*y\) for large real \(s>0\).
Take a zero leading row and column in \(R\) and factor the
remaining positive-semidefinite block by induction.
Uniqueness at a positive pivot follows from the same row
comparison as before. At a zero pivot the normalization
forces the entire first row of \(R\) to be zero. Without this
normalization, \(\operatorname{diag}(0,1)\) also has
\(R=\left(\begin{smallmatrix}0&1\\0&0\end{smallmatrix}\right)\)
as an upper triangular factor with nonnegative diagonal.
\end{proof}

\begin{theorem}[Conditional backward error of real Cholesky]
\label{p206:W26}
Use a real symmetric input matrix and compute the lower triangular
Cholesky factor column by column. To form the residual for
\((i,j)\), \(i\ge j\), start with \(a_{ij}\) and apply
\eqref{p206:subtracttrace}
to products \(\hat l_{it}\hat l_{jt}\) for \(t<j\).
Apply a rounded square root to the resulting diagonal residual,
and divide the off-diagonal residuals by that stored diagonal.
Suppose the square roots have positive arguments, the divisions
have nonzero divisors, and the relative model holds for every
operation, including the square roots. If \((3n+3)u<1\), then
\[
        A+\Delta A=\hat L\hat L^*,\qquad
        |\Delta A|\le\gamma_{3n+3}|\hat L|\,|\hat L^*|.
\]
\end{theorem}

\begin{proof}
For \(i>j\) the terminal residual equals
\(\hat l_{ij}\hat l_{jj}/(1+\delta_{ij})\). For \(i=j\)
it equals \(\hat l_{jj}^2/(1+\delta_{jj})^2\), because the
rounded square root of that residual is \(\hat l_{jj}\).
The terminal product in \eqref{p206:subtracttrace} thus has at most \(j+1\)
signed factors, and every earlier product has at most \(j-1\).
Apply Lemma~\ref{p206:W06} with these counts and use
\(j+1\le 3n+3\). The lower-triangular estimates imply the
upper-triangular ones by symmetry. Define the correction
as \(\hat L\hat L^*-A\).
The hypotheses separate successful rounded execution from
the positive-definite exact factorization. For example, use binary
arithmetic with precision four and round to nearest on
\[
 A=\begin{pmatrix}9/32&1/4\\1/4&1/4\end{pmatrix}.
\]
The determinant is \(1/128>0\) and \(a_{11}>0\), so \(A\) is
positive definite. The exact first square root is between
\(1/2\) and \(17/32\), the midpoint of \(1/2\) and its
successor \(9/16\), since
\((17/32)^2=289/1024>9/32\). It rounds to \(1/2\). The first
multiplier is then exactly \(1/2\), and its rounded square is
exactly \(1/4\). The computed second diagonal residual is zero,
even though the exact Schur complement is positive.
\end{proof}

For the exact factors and conditional Cholesky bounds see
\cite[Chapter~10]{higham}, \cite[Section~4.2]{golub}, and
\cite[entry 14.qr\_cholesky]{d05}.

\section{Exact QR and least squares}

\begin{theorem}[Reduced QR with a fixed diagonal convention]
\label{p206:W27}
For \(A\in\mathbb K^{m\times n}\) of full column rank, with
\(\mathbb K=\mathbb R\) or \(\mathbb C\), there are unique
\(Q,R\) such that \(A=QR\), \(Q^*Q=I_n\), and \(R\) is upper
triangular with positive real diagonal.
\end{theorem}

\begin{proof}
Process the columns \(a_j\) in order. Remove their projections
onto the previous orthonormal \(q_i\), obtaining
\(v_j=a_j-\sum_{i<j}q_i(q_i^*a_j)\). This vector is nonzero
because otherwise \(a_j\) is in the span of the preceding
columns. Put \(r_{jj}=\|v_j\|_2\), \(q_j=v_j/r_{jj}\),
and \(r_{ij}=q_i^*a_j\) for \(i<j\). The projection construction
gives orthonormal columns and \(A=QR\). The empty factorization
fits the same statement. For uniqueness, \(R^*R=A^*A\) and
\(A^*A\) is positive definite when \(n>0\), so \(R\) is
determined by Theorem~\ref{p206:W24}. Then \(Q=AR^{-1}\).
\end{proof}

\begin{theorem}[Least squares and the normal equations]
\label{p206:W28}
For a linear map \(A\) between finite-dimensional real or complex
inner-product spaces and a vector \(b\) in the target, a vector \(x\)
minimizes \(\|Ax-b\|_2\) exactly when \(A^*(Ax-b)=0\).
Minimizers always exist. If \(A\) is injective, the minimizer
is unique and equals \((A^*A)^{-1}A^*b\).
\end{theorem}

\begin{proof}
Decompose \(b=b_{\parallel}+b_\perp\) into its orthogonal
components in \(\operatorname{ran}(A)\) and its orthogonal
complement. Then
\[
      \|Ax-b\|_2^2=\|Ax-b_{\parallel}\|_2^2+\|b_\perp\|_2^2 .
\]
The first term can be made zero, and it vanishes exactly when
\(Ax-b\) is orthogonal to the range. The adjoint identifies
this condition with the normal equations. Injectivity implies
\(\langle A^*Av,v\rangle=\|Av\|_2^2>0\) for each \(v\ne0\),
so \(A^*A\) is invertible in finite dimensions and the
stated formula follows.
\end{proof}

\begin{corollary}[Reduced QR solution formula]\label{p206:W29}
If \(A=QR\) is the full-column-rank reduced QR factorization, the
unique least-squares solution is \(x=R^{-1}Q^*b\).
\end{corollary}

\begin{proof}
Substitution into the normal equations gives
\(R^*Rx=R^*Q^*b\). Both triangular factors \(R,R^*\) are
invertible. Cancel the first factor and solve the remaining
triangular equation. Equivalently \(Ax=QQ^*b\) is the
orthogonal projection of \(b\) onto the column space.
\end{proof}

\begin{theorem}[A real Householder reflector]\label{p206:W30}
For \(x\in\mathbb R^d\setminus\{0\}\), set
\[
 \alpha=-\operatorname{sign}_{+}(x_1)\|x\|_2,\qquad
 v=x-\alpha e_1,\qquad
 H(v)=I-\frac{2vv^*}{v^*v},
\]
where \(\operatorname{sign}_{+}(0)=1\). Then
\(|v_1|=|x_1|+\|x\|_2\), \(v\ne0\), \(H(v)\) is orthogonal,
and \(H(v)x=\alpha e_1\).
\end{theorem}

\begin{proof}
The sign choice proves the identity for \(|v_1|\). The rank-one
matrix \(vv^*/(v^*v)\) is the self-adjoint projection onto
\(\mathbb Rv\). Squaring \(I\) minus twice that projection
gives \(I\), so \(H(v)\) is orthogonal. Because
\(\alpha^2=\|x\|_2^2\),
\[
     v^*v=2\|x\|_2^2-2\alpha x_1=2v^*x .
\]
The formula for \(H(v)x\) now gives \(x-v=\alpha e_1\).
The zero vector may be left fixed by using the identity.
\end{proof}

\begin{theorem}[Exact Householder QR and the condition number of \(A^*A\)]
\label{p206:W31}
For a real \(m\)-by-\(n\) matrix with \(m\ge n\), successive
Householder reflections applied to the column tails construct a
factorization \(A=QR\) with \(Q^*Q=I_n\) and \(R\) upper
triangular. For full column rank one can choose the signs to
give positive diagonal. If \(n>0\) and \(A\) has full column
rank, then
\[
       \kappa_2(A^*A)=\kappa_2(A)^2,\qquad
       \kappa_2(A)=\frac{\sigma_{\max}(A)}{\sigma_{\min}(A)}.
\]
Here \(\sigma_{\min}(A)>0\) is the least singular value on
the domain.
\end{theorem}

\begin{proof}
At step \(j\), apply Theorem~\ref{p206:W30} on coordinates
\(j,\ldots,m\) of column \(j\), and extend the reflector by
the identity on the earlier coordinates. An earlier column
has zero tail on this block and is unchanged. A zero current
tail is also left unchanged. After all steps the matrix is
\(\binom{R}{0}\). Inverting the product of the orthogonal
updates gives the full orthogonal factor, whose first \(n\)
columns give \(Q\). For full column rank \(R\) is invertible,
so its diagonal entries are nonzero. Multiplying the
corresponding columns of \(Q\) and rows of \(R\) by the same
sign makes each diagonal positive and preserves \(QR\).
The spectral theorem applied to \(A^*A\) gives eigenvalues
\(\sigma_j(A)^2\). Its norm and inverse norm are the largest
of these and the reciprocal of the least, respectively. This
gives the squared condition number.
\end{proof}

For the equivalent projection and factorization interfaces see
\cite[entries 14.least\_squares, 14.normal\_equations, and
14.qr\_cholesky]{d05}, \cite[Sections~5.1--5.3]{golub}, and
\cite[Chapters~19--20]{higham}.

\section{Rounded QR and least-squares sensitivity}

\begin{theorem}[Columnwise backward error of a real Householder trace]
\label{p206:W32}
Assume \(m\ge n\). At step \(j\) of the procedure in
Theorem~\ref{p206:W31}, let \(d=m-j+1\) be the active row count
and \(x\) the stored tail of the pivot column. If \(x=0\), do
nothing. Otherwise compute a sequential dot product of \(x\) with
itself, take its rounded square root \(\hat\nu\), set
\(\hat\alpha=-\operatorname{sign}_+(x_1)\hat\nu\), and store
\(w=(\operatorname{fl}(x_1-\hat\alpha),x_2,\ldots,x_d)\).
Replace the active pivot tail by \(\hat\alpha e_1\).
For each trailing stored tail \(y\), compute sequential dot
products \(\hat p\) for \(w^*y\) and \(\hat q\) for \(w^*w\), and use
\[
 \hat t=\operatorname{fl}\bigl(\operatorname{fl}(2\hat p)/\hat q\bigr),
 \qquad
 \hat y_i=\operatorname{fl}\bigl(y_i-\operatorname{fl}(\hat t w_i)\bigr).
\]
Suppose all listed operations are safe, signs and copying are exact,
and \(du<1\) and \(\gamma_d<1\) for every nontrivial step.
There is a product of exact orthogonal reflectors whose first \(n\)
columns form \(Q\), and a matrix \(E\) with \(A+E=Q\hat R\)
for the stored \(\hat R\), such that
\[
       \|E_{:i}\|_2
           \le \left(\prod_{j=1}^{n}(1+c_j)-1\right)\|A_{:i}\|_2 .
\]
The constants \(c_j\) are given explicitly below in terms of
\(d\) and \(u\), with \(c_j=0\) for a skipped step.
\end{theorem}

\begin{proof}
Fix a nontrivial step, write \(\gamma=\gamma_d\), and let
\(\nu=\|x\|_2>0\). The dot-product bound
(Theorem~\ref{p206:W07}) places the computed squared norm in
\([(1-\gamma)\nu^2,(1+\gamma)\nu^2]\). It is positive.
The relative square-root error then yields
\[
 |\hat\nu-\nu|\le a\nu,\qquad
 a=\max\{1-(1-u)\sqrt{1-\gamma},\,
                     (1+u)\sqrt{1+\gamma}-1\}.
\]
With \(\alpha,v\) as in Theorem~\ref{p206:W30}, only the first
coordinate of \(w\) differs from the corresponding coordinate
of \(v\). Its error is bounded by
\[
 \|w-v\|_2\le a\nu+u(|x_1|+\hat\nu)\le b_0\nu,
 \qquad b_0=a+u(2+a).
\]
The sign choice makes the unrounded first component have
magnitude \(|x_1|+\hat\nu>0\); multiplication by its positive
relative factor preserves nonvanishing. Thus \(w\ne0\).
For nonzero \(v,w\), let \(p=v/\|v\|_2\) and \(q=w/\|w\|_2\).
Adding and subtracting \(w/\|v\|_2\) gives
\(\|p-q\|_2\le 2\|v-w\|_2/\|v\|_2\).
The identity \(pp^*-qq^*=(p-q)p^*+q(p-q)^*\) then gives
\[
   \|H(v)-H(w)\|_2\le4\|p-q\|_2\le8b_0,
\]
using \(\|v\|_2\ge\nu\). In particular the stored pivot replacement
differs from \(H(w)x\) by at most \((a+8b_0)\|x\|_2\).

For a trailing tail \(y\), write \(p_0=w^*y\),
\(q_0=\|w\|_2^2\), and \(t_0=2p_0/q_0\). The dot-product
bounds give
\[
 |\hat p-p_0|\le\gamma\|w\|_2\|y\|_2,\qquad
       \hat q=q_0(1+\varepsilon),\quad |\varepsilon|\le\gamma .
\]
Here \(\sum_i|w_i y_i|\le\|w\|_2\|y\|_2\) and \(q_0>0\).
The scale and divide operations give
\(\hat t=(2\hat p/q_0)r\) with
\(r=(1+\delta_1)(1+\delta_2)/(1+\varepsilon)\) and
\(|\delta_k|\le u\). Put
\[
 R_0=\frac{(1+u)^2}{1-\gamma},\qquad
 D_0=\max\left\{R_0-1,\,1-\frac{(1-u)^2}{1+\gamma}\right\}.
\]
Then \(0<r\le R_0\) and \(|r-1|\le D_0\). Expanding
\(r\hat p-p_0=(r-1)p_0+r(\hat p-p_0)\) and using Cauchy--Schwarz
gives
\[
 \begin{aligned}
  |\hat t-t_0|\|w\|_2&\le 2(D_0+R_0\gamma)\|y\|_2,\\
  |\hat t|\|w\|_2&\le 2R_0(1+\gamma)\|y\|_2 .
 \end{aligned}
\]
In each component the product and subtraction relative errors
give
\[
 |\hat y_i-(y_i-\hat t w_i)|
      \le u|y_i|+(2u+u^2)|\hat t w_i|.
\]
This follows by expanding the subtraction as
\((y_i-\hat t w_i(1+\delta))(1+\epsilon)\) and bounding the two
coefficients. Since \(H(w)y=y-t_0w\), the trailing-column
error is bounded by \(c_{\rm tail}\|y\|_2\), where
\[
   c_{\rm tail}
     =2(D_0+R_0\gamma)+u+2(2u+u^2)R_0(1+\gamma).
\]
Take \(c_j=\max\{a+8b_0,c_{\rm tail}\}\) for this step.

Let \(B_j\) be the full stored matrix after \(j\) updates and
\(H_j\) the exact \(H(w)\) extended by the identity on the
earlier rows. For each column, the preceding estimates say
\(B_j=H_jB_{j-1}+F_j\) with
\(\|(F_j)_{:i}\|_2\le c_j\|(B_{j-1})_{:i}\|_2\).
Previously cleared columns stay exactly unchanged since their
tails are zero. The full column norm dominates its tail norm,
so the same inequality holds for the full matrices.
Orthogonality gives
\(\|(B_j)_{:i}\|_2\le\prod_{k\le j}(1+c_k)\|A_{:i}\|_2\).
Pull the final matrix \(\binom{\hat R}{0}\) back by the
inverse orthogonal product. Its difference from \(A\) is the
sum of the local corrections transformed by orthogonal factors.
For column \(i\) the bound is
\[
   \sum_{j=1}^n c_j\prod_{k<j}(1+c_k)\|A_{:i}\|_2
        =\left(\prod_{j=1}^n(1+c_j)-1\right)\|A_{:i}\|_2 .
\]
This constructs the exact implicit \(Q\) claimed in the theorem.
If \(n=0\) the bound is zero. If an explicit \(Q\) is formed
by further rounded operations, its error or orthogonality
defect must be bounded separately.
\end{proof}

\begin{theorem}[QR residual and orthogonality certificate]
\label{p206:W33}
Let \(m\ge n\), let \(A,\hat Q\) be \(m\)-by-\(n\) matrices
and \(\hat R\) be \(n\)-square, and suppose
\[
       \|A-\hat Q\hat R\|_2\le\rho,\qquad
       \|\hat Q^*\hat Q-I\|_2\le\eta<1,\qquad \eta\ge0.
\]
There is an \(m\)-by-\(n\) matrix \(\widetilde Q\) with
\(\widetilde Q^*\widetilde Q=I_n\) and
\[
 \|A-\widetilde Q\hat R\|_2
     \le \rho+\frac{\eta}{1+\sqrt{1-\eta}}\,\|\hat R\|_2 .
\]
\end{theorem}

\begin{proof}
Put \(G=\hat Q^*\hat Q\). Its eigenvalues lie in
\([1-\eta,1+\eta]\) by the Rayleigh quotient bound for \(G-I\),
so its positive square root and inverse square root exist by
orthogonal or unitary diagonalization. Define
\(\widetilde Q=\hat QG^{-1/2}\). Its columns are orthonormal, and
\[
 \hat Q-\widetilde Q=\widetilde Q(G^{1/2}-I),\qquad
 \|\hat Q-\widetilde Q\|_2=\|G^{1/2}-I\|_2.
\]
Multiplication by \(\widetilde Q\) preserves vector norms.
For \(\lambda\in[1-\eta,1+\eta]\),
\[
       |\sqrt{\lambda}-1|
       =\frac{|\lambda-1|}{1+\sqrt{\lambda}}
       \le\frac{\eta}{1+\sqrt{1-\eta}} .
\]
The spectral norm bound follows, and the triangle inequality
and submultiplicativity applied to
\(A-\widetilde Q\hat R=(A-\hat Q\hat R)+(\hat Q-\widetilde Q)\hat R\)
complete the proof. The empty-column case is vacuous.
This certificate uses supplied bounds \(\rho,\eta\). Their
verification for any explicit stored \(\hat Q,\hat R\) is a
separate part of a numerical argument.
\end{proof}

\begin{theorem}[Derivative of a full-column-rank least-squares solution]
\label{p206:W34}
For real \(A\) of full column rank and a real vector \(b\),
let \(x=(A^*A)^{-1}A^*b\), \(r=b-Ax\), and
\(A^\dagger=(A^*A)^{-1}A^*\). As a function of the matrix and
right-hand side, the Fr\'echet derivative of \(x\) in the
direction \((E,f)\) is
\[
             A^\dagger(f-Ex)+(A^*A)^{-1}E^*r .
\]
\end{theorem}

\begin{proof}
Put \(B=A^*A\), \(B_E=(A+E)^*(A+E)\). The difference
\(B_E-B=A^*E+E^*A+E^*E\) is \(O(\|E\|_2)\) as \(E\to0\).
Theorem~\ref{p206:W10} makes \(B_E\) invertible near \(E=0\)
with uniformly bounded inverse. The inverse identity
\(B_E^{-1}-B^{-1}=-B_E^{-1}(B_E-B)B^{-1}\)
gives \(B_E^{-1}-B^{-1}=O(\|E\|_2)\). In particular \(A+E\)
has full column rank there. Subtract its normal equations
from the equation at the old \(x\), in the form
\[
 \begin{aligned}
  B_E(x(A+E,b+f)-x)
      &=(A+E)^*(b+f-(A+E)x)\\
      &= A^*(f-Ex)+E^*r+E^*(f-Ex).
 \end{aligned}
\]
In the last step \(A^*r=0\). The last term is
\(O((\|E\|_2+\|f\|_2)^2)\), and replacing \(B_E^{-1}\) by
\(B^{-1}\) on the linear terms makes another error of this
order. This proves Fr\'echet differentiability with the stated
derivative. The residual-dependent term vanishes for a consistent
system and can affect sensitivity for inconsistent data.
\end{proof}

\begin{theorem}[Rank-deficient least squares and the minimum-norm solution]
\label{p206:W35}
For an arbitrary real or complex matrix \(A\), the set of
least-squares minimizers is \(A^\dagger b+\ker A\), with
\(A^\dagger b\) the unique minimizer of least Euclidean norm.
\end{theorem}

\begin{proof}
Use the exact singular value decomposition
\cite[entry 14.svd]{d05}. For the positive singular values
\(\sigma_i\), write \(Av_i=\sigma_i u_i\) on the orthonormal
right singular vectors and let \(A^\dagger b=
\sum_{\sigma_i>0}\sigma_i^{-1}v_i(u_i^*b)\).
The vectors \(u_i\) span \(\operatorname{ran}(A)\), while the
vectors \(v_i\) span \((\ker A)^\perp\).
Consequently \(AA^\dagger b\) is the projection of \(b\)
onto the range. The proof of Theorem~\ref{p206:W28} identifies
all minimizers as the inverse image of that projection, namely
\(A^\dagger b+\ker A\). Orthogonality of \(A^\dagger b\) to
the kernel makes the squared norm of \(A^\dagger b+k\)
equal to \(\|A^\dagger b\|_2^2+\|k\|_2^2\).
This also handles the rank-zero case with an empty sum.
\end{proof}

For Householder error analysis and least-squares perturbations,
compare \cite[Chapters~19--20]{higham}. The exact rank-deficient
formula also appears in \cite[Section~5.5]{golub} and
\cite[entries 14.least\_squares and 14.svd\_pseudoinverse]{d05}.
\chapter{Spectral Perturbation and Low-Rank Approximation}

We work over \(\mathbb C\), with the real versions obtained when the
data and the asserted decompositions are real. The Euclidean
operator norm is \(\|\cdot\|_2\) and the Frobenius norm is
\(\|\cdot\|_F\). We use the orthonormal spectral theorem for
Hermitian and normal matrices and the finite singular value
decomposition, as proved in \cite[Modules 13 and 14]{d05}.
A rectangular matrix \(A:\mathbb K^n\to\mathbb K^m\) has an
orthonormal basis of right singular vectors, obtained from
\(A^*A\). The singular values are ordered nonincreasingly and
are padded with zeros beyond the positive rank. A Hermitian
matrix of size \(n\) has eigenvalues in nondecreasing order.

\section{Hermitian eigenvalues and residuals}

\begin{theorem}[Courant--Fischer]\label{p206:W36}
For Hermitian \(A\) of size \(n\) write \(R_A(x)=x^*Ax\)
on unit vectors. For \(1\le i\le n\),
\[
 \lambda_i(A)
  = \min_{\dim S=i}\ \max_{\substack{x\in S\\\|x\|_2=1}}R_A(x)
  = \max_{\dim S=n-i+1}\ \min_{\substack{x\in S\\\|x\|_2=1}}R_A(x).
\]
\end{theorem}
\begin{proof}
Choose an orthonormal eigenbasis \(v_1,\ldots,v_n\) in that order.
On the span of the first \(i\) basis vectors the maximum Rayleigh
quotient is \(\lambda_i\). Every \(i\)-dimensional subspace
intersects \(\operatorname{span}(v_i,\ldots,v_n)\) nontrivially
because the dimensions sum to \(n+1\). A unit vector in that
intersection has Rayleigh quotient at least \(\lambda_i\).
This proves the first formula. Interchanging the two spans and
inequalities gives the second. The extreme values on a subspace
are attained by compactness of its unit sphere. Repeated eigenvalues
require no choice of a particular eigenbasis within their spaces.
\end{proof}

\begin{theorem}[Weyl's Hermitian perturbation bound]\label{p206:W37}
If \(A,E\) are Hermitian of the same size, then for every
eigenvalue index,
\[
                |\lambda_i(A+E)-\lambda_i(A)|\le\|E\|_2 .
\]
\end{theorem}
\begin{proof}
On a unit vector, the Rayleigh quotients differ by
\(x^*Ex\), of modulus at most \(\|E\|_2\) by Cauchy--Schwarz.
The same bound therefore holds after taking the maximum in each
fixed subspace and then the minimum over all subspaces of
dimension \(i\) in Theorem~\ref{p206:W36}.
\end{proof}

\begin{theorem}[Gershgorin and a Hermitian residual certificate]
\label{p206:W38}
Every eigenvalue of a complex square matrix belongs to one of the
row disks with centers \(a_{ii}\) and radii
\(\sum_{j\ne i}|a_{ij}|\). The same conclusion holds for column
disks. If \(A\) is Hermitian, \(v\) is a unit vector and
\(\mu\in\mathbb R\), then
\[
 \operatorname{dist}(\mu,\sigma(A))\le\|(A-\mu I)v\|_2 .
\]
For a group of eigenvalues, let \(P\) be the orthogonal projection
onto their joint eigenspaces. If the other eigenvalues all have
distance at least \(\delta>0\) from \(\mu\), then
\(\|(I-P)v\|_2\le\|(A-\mu I)v\|_2/\delta\).
\end{theorem}
\begin{proof}
In an eigenvector choose a coordinate \(i\) of maximum modulus.
It is nonzero, and the \(i\)-th row equation gives
\(|\lambda-a_{ii}|\le\sum_{j\ne i}|a_{ij}|\) after division by
that modulus. The transposed matrix has the same characteristic
polynomial and its row disks are the column disks of \(A\).
For the Hermitian assertions write \(v=\sum_j c_j v_j\) in an
orthonormal eigenbasis. Then
\[
 \|(A-\mu I)v\|_2^2=\sum_j|\lambda_j-\mu|^2|c_j|^2,
 \qquad\sum_j|c_j|^2=1.
\]
Use the minimum spectral distance in the whole sum for the first
bound, and the gap \(\delta\) in the sum over the complement for
the second.
\end{proof}
These comparisons are developed in \cite[Sections 7.2 and 8.1--8.2]{golub}.
The row disk theorem also has a fixed Mathlib and MathGenome
implementation \cite{mathlib,d05}.

\section{Diagonalizable perturbations}

\begin{theorem}[Bauer--Fike with a specified subordinate norm]
\label{p206:W39}
Let \(A=V\Lambda V^{-1}\) with \(\Lambda\) diagonal, and let
\(\mu\) be an eigenvalue of \(A+E\).
For any subordinate norm and \(\mu\notin\sigma(A)\),
\[
 1\le\|(\mu I-\Lambda)^{-1}\|\ \|V^{-1}EV\|.
\]
If the norm of a diagonal matrix is the maximum modulus of its
diagonal entries, then
\[
      \min_i|\mu-\lambda_i|\le\|V\|\,\|V^{-1}\|\,\|E\|.
\]
This hypothesis holds for the subordinate \(\ell_p\) norms,
\(1\le p\le\infty\).
\end{theorem}
\begin{proof}
For a nonzero eigenvector \(x\) put \(z=V^{-1}x\).
If \(\mu\) is outside the spectrum of \(A\), its equation
gives \(z=(\mu I-\Lambda)^{-1}V^{-1}EVz\).
Take norms and divide by \(\|z\|>0\).
For an \(\ell_p\) norm the diagonal operator is bounded above
by the maximum diagonal modulus directly from the vector norm,
and a coordinate vector attains it. The stated diagonal norm
property turns the inverse norm into
\(1/\min_i|\mu-\lambda_i|\). Submultiplicativity finishes the
estimate. When \(\mu\) is already an eigenvalue of \(A\),
the left side of the spectral-distance assertion is zero.
\end{proof}
The subordinate \(p\)-norm statement is \cite[Theorem 7.2.2]{golub}.
An arbitrary subordinate norm need not have the diagonal property.
For a concrete failure, take
\[
 A=\begin{pmatrix}0&0\\0&1\end{pmatrix},\quad
 S=\begin{pmatrix}1&M\\0&1\end{pmatrix},\quad
 SES^{-1}=\begin{pmatrix}0&0\\1/M&0\end{pmatrix},\qquad M>2.
\]
Use the vector norm \(\|x\|=\|Sx\|_\infty\). Its subordinate
norm on \(E\) is \(1/M\). The diagonalization of \(A\)
may use \(V=I\), whose condition number is one in this norm.
But \(S(A+E)S^{-1}\) has rows \((0,M)\) and \((1/M,1)\).
Its eigenvalues are \((1\pm\sqrt5)/2\), both at distance
\((\sqrt5-1)/2\) from \(\{0,1\}\).
This exceeds \(1/M\). The resolvent inequality still holds.

\section{Singular values and optimal approximation}

\begin{theorem}[Singular-value perturbations]\label{p206:W40}
For same-size real or complex matrices \(A,A+E\), with
singular values ordered and padded as above,
\[
                |\sigma_i(A+E)-\sigma_i(A)|\le\|E\|_2 .
\]
\end{theorem}
\begin{proof}
On the \(n\)-dimensional right-hand space of \(A\), for
\(1\le i\le n\),
\[
 \sigma_i(A)
     =\min_{\dim S=n-i+1}\ \max_{\substack{x\in S\\\|x\|_2=1}}\|Ax\|_2 .
\]
The span of the last \(n-i+1\) right singular vectors gives
the upper bound. Every subspace of that dimension intersects
the span of the first \(i\) right singular vectors, and a unit
vector in the intersection gives the lower bound. This uses
the SVD formula for \(\|Ax\|_2^2\), including all its zero
directions. Since \(\|Ax\|_2\) and \(\|(A+E)x\|_2\) differ
by at most \(\|E\|_2\) on any unit vector, compare the
two subspace extrema. Indices beyond \(n\) contribute zero;
when the right space has dimension zero the result is immediate.
\end{proof}

\begin{theorem}[Best rank-\(k\) approximation in operator norm]
\label{p206:W41}
For an integer \(k\ge0\), truncating an SVD after its first
\(k\) singular terms gives a minimizer over all matrices of
rank at most \(k\). The minimum of \(\|A-B\|_2\) is
\(\sigma_{k+1}(A)\).
\end{theorem}
\begin{proof}
The residual of the truncated SVD acts on the remaining singular
directions with largest coefficient \(\sigma_{k+1}\), so its
operator norm is that coefficient. When \(k\ge\operatorname{rank}A\),
the truncation is \(A\) and the minimum is zero.
Otherwise \(\ker B\) has dimension at least \(n-k\) for any
competitor \(B\), and intersects the span of the first \(k+1\)
right singular vectors. On a unit vector \(x\) there,
\((A-B)x=Ax\) has norm at least \(\sigma_{k+1}\).
The lower bound follows.
\end{proof}

\begin{theorem}[Best rank-\(k\) approximation in Frobenius norm]
\label{p206:W42}
The same SVD truncation minimizes \(\|A-B\|_F\) over matrices
of rank at most \(k\), with squared minimum
\(\sum_{i>k}\sigma_i(A)^2\).
\end{theorem}
\begin{proof}
For \(k\ge\operatorname{rank} A\) the claim holds by taking \(B=A\).
For \(k<\operatorname{rank} A\), choose an orthonormal basis
\(y_1,\ldots,y_{n-k}\) of an \((n-k)\)-dimensional subspace of
\(\ker B\). Extend it to an orthonormal basis of the right space.
By the Frobenius sum formula,
\[
 \|A-B\|_F^2
    \ge\sum_{j=1}^{n-k}\|Ay_j\|_2^2
    =\sum_{i=1}^n\sigma_i^2 w_i,\qquad
       w_i=\sum_{j=1}^{n-k}|v_i^*y_j|^2 .
\]
Here \(v_i\) are the right singular vectors. By orthogonal
projection, \(0\le w_i\le1\), and summing the Parseval identity
for each \(y_j\) gives \(\sum_i w_i=n-k\).
The difference from \(\sum_{i>k}\sigma_i^2\) is
\(\sum_{i\le k}\sigma_i^2w_i-\sum_{i>k}\sigma_i^2(1-w_i)\).
The weights in the two sums have equal total, and the first
set of coefficients is at least the second set, so the difference
is nonnegative. For \(k=0\), all weights equal one.
The residual of the truncated SVD attains the claimed value.
\end{proof}
The SVD constructions and norm identities used here are in
\cite[Section 14]{d05}; see also \cite[Section 2.4]{golub}.

\begin{theorem}[Hoffman--Wielandt and a Frobenius singular-value bound]
\label{p206:W43}
For normal \(n\times n\) matrices \(A,B\) with eigenvalues
\(\lambda_i,\mu_i\), there is a permutation \(\pi\) with
\[
        \sum_i|\lambda_i-\mu_{\pi(i)}|^2\le\|A-B\|_F^2 .
\]
For Hermitian matrices one may match the eigenvalues in increasing
order. For rectangular matrices \(A,B\) of the same size,
\[
        \sum_i\bigl(\sigma_i(A)-\sigma_i(B)\bigr)^2\le\|A-B\|_F^2 .
\]
\end{theorem}
\begin{proof}
Use orthonormal eigenbases \(u_i,v_j\) for the normal matrices.
The mixed-basis entry of \(A-B\) is
\(u_i^*(A-B)v_j=(\lambda_i-\mu_j)u_i^*v_j\).
Indeed \(u_i^*A=\lambda_i u_i^*\) by the orthonormal
diagonalization, and \(Bv_j=\mu_jv_j\).
Frobenius norm is unchanged by orthonormal changes on either side,
so
\[
  \|A-B\|_F^2=\sum_{i,j}|\lambda_i-\mu_j|^2\,|u_i^*v_j|^2.
\]
The matrix \(P_{ij}=|u_i^*v_j|^2\) has row and column sums one
by Parseval. The Birkhoff--von Neumann decomposition proved
in \cite[Chapter 8, W42]{p205} writes \(P\) as a convex
combination of permutation matrices. The displayed cost is
linear in \(P\), so one permutation has cost at most its average.
For Hermitian matrices the eigenvalues are real. If
\(a\le b\) and \(c\le d\), swapping the crossed pairing
\((a,d),(b,c)\) to \((a,c),(b,d)\) decreases the squared cost by
\(2(b-a)(d-c)\ge0\). Repeated removal of inversions shows that
the sorted pairing is optimal, including ties.

For rectangular \(A\), use the Hermitian matrix on the direct sum
of the left and right spaces given by
\[
                  H(A)=\begin{pmatrix}0&A\\A^*&0\end{pmatrix}.
\]
For each positive singular value with unit pair \(u,v\),
the vectors \((u,v)\) and \((u,-v)\) have eigenvalues
\(\sigma\) and \(-\sigma\); the remaining kernel directions give
zero. More explicitly, with \(q=\min(m,n)\) the increasing list
is \(-\sigma_1,\ldots,-\sigma_q\), followed by \(|m-n|\) zeros,
followed by \(\sigma_q,\ldots,\sigma_1\).
The paired squared spectral differences for \(H(A),H(B)\)
therefore sum to twice the sum of squared singular-value
differences. Meanwhile
\(\|H(A)-H(B)\|_F^2=2\|A-B\|_F^2\).
The Hermitian result gives the claim.
\end{proof}
See \cite{hoffman,mirsky} for the classical spectral and
singular-value inequalities.

\pagebreak
\section{Power and inverse iterations}

\begin{theorem}[The dominant direction of power iteration]
\label{p206:W44}
Let \(A\) be diagonalizable and let its nonzero eigenvalue
\(\lambda_1\) have modulus strictly larger than that of every
other distinct eigenvalue. Let \(p_1\) be the spectral component
of \(x_0\) in the \(\lambda_1\) eigenspace and assume \(p_1\ne0\).
Then the unit directions \(z_j=A^jx_0/\|A^jx_0\|_2\),
after multiplication by \((|\lambda_1|/\lambda_1)^j\),
converge to \(p_1/\|p_1\|_2\) with error \(O(r^j)\), where
\(r=\max_{\lambda\ne\lambda_1}|\lambda/\lambda_1|<1\)
and the maximum of the empty set is zero.
\end{theorem}
\begin{proof}
Decompose \(x_0=p_1+\sum_{\lambda\ne\lambda_1}p_\lambda\)
in the direct sum of eigenspaces. Since the \(\lambda_1\) component
of \(A^jx_0\) is nonzero, all directions are defined.
Put \(y_j=\lambda_1^{-j}A^jx_0\). For \(j\ge1\),
\(\|y_j-p_1\|_2\le r^j\sum_{\lambda\ne\lambda_1}\|p_\lambda\|_2\).
For any nonzero \(y,p\) one has
\[
 \left\|\frac{y}{\|y\|_2}-\frac{p}{\|p\|_2}\right\|_2
  \le \frac{\|y-p\|_2+|\|y\|_2-\|p\|_2|}{\|p\|_2}
  \le \frac{2\|y-p\|_2}{\|p\|_2}.
\]
Use this with \(y_j,p_1\); its normalized \(y_j\) is the stated
phase correction. If all other eigenvalues are zero, the first
iterate has the exact corrected direction. In the real case the
phase correction is a sign. Multiplicity of the dominant
eigenvalue within its eigenspace does not affect this proof.
\end{proof}

\begin{corollary}[Fixed-shift inverse iteration]\label{p206:W45}
Let \(A\) be diagonalizable and \(\mu\notin\sigma(A)\).
Suppose one distinct eigenvalue \(\lambda_1\) uniquely minimizes
\(|\lambda-\mu|\), with a nonzero initial component in its
eigenspace. The exact iteration that solves
\((A-\mu I)y=x\) and normalizes \(y\) converges in direction
to that component, with the phase convention of
Theorem~\ref{p206:W44} and geometric factor
\[
       \max_{\lambda\ne\lambda_1}
                    \frac{|\lambda_1-\mu|}{|\lambda-\mu|}.
\]
With a single distinct eigenvalue the directions are already exact
up to phase.
\end{corollary}
\begin{proof}
The operator \((A-\mu I)^{-1}\) has the same eigenspaces and
eigenvalues \((\lambda-\mu)^{-1}\). The closest eigenvalue
of \(A\) gives the uniquely largest modulus for this operator.
Apply the power iteration result. Normalizing after each step
gives the same direction as normalizing its power applied to
the original vector. The proof concerns exact solves.
A residual of a rounded solve can be converted to a forward error
using Theorem~\ref{p206:W14}, with the inverse norm of the shifted
matrix.
\end{proof}
\chapter{Polynomial Interpolation and Approximation}
\label{p206:ch06}

Let \(\mathcal P_n(K)\) denote the polynomials over \(K\) of degree at
most \(n\), including the zero polynomial. When \(K=\mathbb R\) and an
interval is specified, we also view these polynomials as functions on
that interval. Distinctness of the nodes is essential for ordinary
interpolation. Repeated nodes will encode higher-order data.

\section{Lagrange, Newton, and barycentric forms}

\begin{proposition}[Lagrange cardinal polynomials]\label{p206:W46}
Let \(K\) be a field and \(x_0,\ldots,x_n\) pairwise distinct elements
of \(K\). For prescribed values \(y_0,\ldots,y_n\), set
\[
 \ell_i(X)=\prod_{j\ne i}\frac{X-x_j}{x_i-x_j},
 \qquad
 p(X)=\sum_{i=0}^n y_i\ell_i(X).
\]
Then \(\ell_i\in\mathcal P_n(K)\), \(\ell_i(x_j)=\delta_{ij}\), and
\(p(x_i)=y_i\) for each \(i\). The singleton case has
\(\ell_0=1\). The empty sum constructs the zero polynomial for empty
data.
\end{proposition}

\begin{proof}
All denominators are nonzero. Each \(\ell_i\) is a product of \(n\)
linear factors scaled by field elements. At \(x_i\) every factor is
\(1\). At \(x_j\), \(j\ne i\), the factor with numerator \(X-x_j\)
vanishes. Evaluation of the sum gives the claim.
\end{proof}

\begin{theorem}[Existence and uniqueness]\label{p206:W47}
Data at \(n+1\) distinct nodes determine a unique polynomial in
\(\mathcal P_n(K)\). The \(\ell_i\) form the basis of this space dual
to evaluation at the nodes.
\end{theorem}

\begin{proof}
Proposition~\ref{p206:W46} gives existence. The difference of two
solutions is divisible by all \(X-x_i\), and hence by their product,
since the linear factors are pairwise coprime. This product has
degree \(n+1\), so the degree bound forces the difference to vanish.
Applying the interpolation formula to arbitrary polynomial data
proves the basis assertion. Empty data have no uniqueness in
\(\mathcal P_n(K)\) for \(n\ge0\).
\end{proof}

These formulas are the ordinary distinct-node interpolation of
\cite[Section 2.1.1]{stoer}. The fixed Mathlib revision
\cite{mathlib} provides the corresponding construction, evaluation
and degree lemmas in \texttt{Mathlib/LinearAlgebra/Lagrange.lean}.

\begin{theorem}[Newton divided differences]\label{p206:W48}
For pairwise distinct nodes in a field, define divided differences
for any ordered sublist by \(f[x_i]=y_i\) and
\[
 f[x_i,\ldots,x_j]
 =\frac{f[x_{i+1},\ldots,x_j]-f[x_i,\ldots,x_{j-1}]}{x_j-x_i}.
\]
The divided difference on a list equals the coefficient of the
highest allowed degree in its interpolation polynomial, and is
invariant under permutations of the node-value pairs. In
particular the interpolant on the original list is
\[
 p(X)=\sum_{k=0}^n c_k\prod_{j=0}^{k-1}(X-x_j),
 \qquad c_k=f[x_0,\ldots,x_k].
\]
To evaluate it at \(x\), start with \(r=c_n\) and replace
\(r\) by \(c_k+(x-x_k)r\) for \(k=n-1,\ldots,0\).
\end{theorem}

\begin{proof}
Write \(p_{i,j}\) for the degree-\(\le j-i\) interpolant on
\(x_i,\ldots,x_j\). For \(i<j\) the polynomial
\[
 \frac{(X-x_i)p_{i+1,j}(X)-(X-x_j)p_{i,j-1}(X)}{x_j-x_i}
\]
takes the prescribed value at both end nodes and at every
interior node. For an interior node both polynomials have the
same data value, and the difference of the two linear factors
is \(x_j-x_i\). Its degree is at most \(j-i\), so it equals
\(p_{i,j}\). Comparing the coefficients of \(X^{j-i}\)
gives the recursion from the respective coefficients of
\(X^{j-i-1}\). The singleton case starts the induction.
Uniqueness makes that coefficient independent of node order.

Next \(p_{0,k}-p_{0,k-1}\) has degree at most \(k\) and
vanishes at the first \(k\) nodes. It is \(c_k\) times their
monic nodal product, by its coefficient of \(X^k\).
Summing these differences gives the representation. The backward
recursion follows by successively factoring out the first
linear factor from the remaining terms.
\end{proof}

\begin{proposition}[Barycentric identities]\label{p206:W49}
For \(n+1\) distinct nodes in a field set
\[
 \omega(X)=\prod_j(X-x_j),\qquad
 w_i=\frac{1}{\prod_{j\ne i}(x_i-x_j)}.
\]
At a non-node \(x\), the exact interpolant satisfies
\[
 p(x)=\omega(x)\sum_{i=0}^n\frac{w_i y_i}{x-x_i}
     =\frac{\displaystyle\sum_{i=0}^n w_i y_i/(x-x_i)}
             {\displaystyle\sum_{i=0}^n w_i/(x-x_i)}.
\]
The denominator of the second expression is nonzero.
At \(x=x_i\), evaluation is the prescribed \(y_i\).
\end{proposition}

\begin{proof}
The first formula follows by factoring \(\ell_i(x)\) as
\(\omega(x)w_i/(x-x_i)\). Apply it to the constant data \(y_i=1\).
Their interpolant is \(1\), by uniqueness. Thus the denominator
in the quotient is \(1/\omega(x)\), which is nonzero, and the
second formula follows. A floating-point analysis additionally
needs error bounds for the computed nodes and weights.
See \cite[Section 2]{berrut} for the barycentric forms.
\end{proof}

\section{Interpolation remainders and repeated nodes}

We use Rolle's theorem in its interval form from
\cite[\texttt{K03.05.rolle}]{d03}. Here \(C^r[a,b]\) uses continuous
derivatives through order \(r\) on the closed interval, with the
usual one-sided interpretation at the ends.

\begin{lemma}[Rolle's theorem with repeated zeros]\label{p206:repeated-rolle}
Let \(g\in C^{N-1}[a,b]\). Suppose it has distinct marked points
\(t_1<\cdots<t_s\) with positive integers \(m_i\), with
\(\sum_i m_i=N\), and that \(g^{(j)}(t_i)=0\) for \(0\le j<m_i\).
Then \(g^{(N-1)}\) vanishes at some point of \([t_1,t_s]\).
\end{lemma}

\begin{proof}
For \(N=1\) use the marked point. For \(N\ge2\), at each \(t_i\)
with \(m_i\ge2\), the derivative \(g'\) has the specified vanishing
of order \(m_i-1\). Also, between each pair \(t_i,t_{i+1}\), Rolle's
theorem supplies an interior zero of \(g'\). The new points
are distinct from one another and from the old ones.
Thus \(g'\) has a list of marked zeros of total multiplicity
\(\sum_i(m_i-1)+(s-1)=N-1\). Induction inside the same
closed hull proves the result.
\end{proof}

\begin{theorem}[Lagrange remainder]\label{p206:W50}
Let \(x_0,\ldots,x_n\) be distinct nodes of a real interval and \(x\)
another point of that interval, possibly a node. Suppose \(f\) is
\(C^{n+1}\) there and let \(p\) interpolate the \(f(x_i)\).
There is a \(\xi\) in the closed convex hull of \(x,x_0,\ldots,x_n\)
such that
\[
 f(x)-p(x)=\frac{f^{(n+1)}(\xi)}{(n+1)!}
                   \prod_{i=0}^n(x-x_i).
\]
\end{theorem}

\begin{proof}
For a node \(x\) the identity is zero on both sides. Otherwise put
\(c=(f(x)-p(x))/\omega(x)\), with \(\omega\) the monic nodal product.
Then \(g=f-p-c\omega\) vanishes at the \(n+1\) nodes and at \(x\).
Apply Lemma~\ref{p206:repeated-rolle} with \(N=n+2\) to this
list of distinct zeros. Since \(p^{(n+1)}=0\) and
\(\omega^{(n+1)}=(n+1)!\), the zero of \(g^{(n+1)}\)
gives \(c=f^{(n+1)}(\xi)/(n+1)!\).
\end{proof}

\begin{theorem}[Hermite interpolation and remainder]\label{p206:W51}
Let \(s\ge1\), let \(x_1,\ldots,x_s\) be pairwise distinct nodes of a field,
and \(m_i\ge1\) integers with \(N=\sum_i m_i\).
For each node prescribe a residue class modulo \((X-x_i)^{m_i}\).
There is a unique representative \(p\) of degree less than \(N\)
meeting all the residues.

For real nodes in an interval and \(f\in C^N\) on that interval,
take the residues to be the Taylor polynomials at \(x_i\) through
degree \(m_i-1\). Then \(p^{(j)}(x_i)=f^{(j)}(x_i)\) for
\(0\le j<m_i\). At each point \(x\) of the interval,
\[
 f(x)-p(x)=\frac{f^{(N)}(\xi)}{N!}
                \prod_{i=1}^s(x-x_i)^{m_i}
\]
for some \(\xi\) in the closed hull of the nodes and \(x\).
\end{theorem}

\begin{proof}
Let \(q_i=(X-x_i)^{m_i}\), \(q=\prod_iq_i\), and
\(Q_i=\prod_{j\ne i}q_j\). Since \(Q_i\) and \(q_i\) are
coprime, the polynomial Euclidean algorithm gives \(u_i,v_i\)
with \(u_iQ_i+v_iq_i=1\). If \(a_i\) represents the requested
class, \(\sum_i a_i u_i Q_i\) has the correct residues.
Its remainder on division by the monic polynomial \(q\)
has degree less than \(N\) and the same residues.
The difference of two solutions of degree less than \(N\)
is divisible by \(q\), so vanishes. This is the polynomial
construction underlying the finite Chinese remainder theorem
\cite[\texttt{K07.09.crt\_finite\_integer}]{d07}.
In arbitrary characteristic the residues are the data.
Over \(\mathbb R\), the coefficient of \((X-x_i)^j\)
is \(p^{(j)}(x_i)/j!\), which proves the equivalence
with the derivative conditions.

At a node the error is zero. Otherwise set
\(c=(f(x)-p(x))/q(x)\) and \(g=f-p-cq\).
At \(x_i\), this has a marked zero of multiplicity \(m_i\),
because \(q\) and its first \(m_i-1\) derivatives vanish
there. There is also the new zero \(x\), giving total
multiplicity \(N+1\). Lemma~\ref{p206:repeated-rolle} gives
\(g^{(N)}(\xi)=0\). As \(p^{(N)}=0\) and \(q^{(N)}=N!\),
the desired identity follows.
\end{proof}

\section{Lebesgue constants and node geometry}

\begin{theorem}[Interpolation norm and near-best error]\label{p206:W52}
Let \(a<b\) and choose \(n+1\) distinct nodes in \([a,b]\).
Let \(I_n:C[a,b]\to\mathcal P_n(\mathbb R)\) be interpolation
with the uniform norm on both spaces. Then
\[
 \lVert I_n\rVert=\Lambda_n
   :=\max_{x\in[a,b]}\sum_{i=0}^n|\ell_i(x)|,
 \qquad
 \lVert f-I_nf\rVert_\infty
 \le(1+\Lambda_n)\inf_{p\in\mathcal P_n(\mathbb R)}
                   \lVert f-p\rVert_\infty .
\]
\end{theorem}

\begin{proof}
The sum defining \(\Lambda_n\) is continuous, so it attains its
maximum at some \(x_*\). The interpolation formula gives the upper
norm bound. Prescribe at each node the sign of \(\ell_i(x_*)\),
choosing \(0\) when \(\ell_i(x_*)=0\). Connect these values by
linear interpolation between adjacent nodes in increasing order,
and extend constantly up to the two interval ends. This gives a
continuous function of uniform norm at most \(1\) with interpolant
equal to \(\Lambda_n\) at \(x_*\). Hence the norm is exactly
\(\Lambda_n\). For any polynomial \(p\in\mathcal P_n(\mathbb R)\),
uniqueness gives \(I_np=p\). Thus \(f-I_nf=(1-I_n)(f-p)\).
Take norms and then the infimum. See \cite[Chapter 15]{trefapprox}.
\end{proof}

\begin{theorem}[Chebyshev nodes and a uniform approximation rate]\label{p206:W53}
For \(m\ge1\), let \(T_m(\cos\theta)=\cos(m\theta)\), and let
\(\theta_j=(j+\tfrac12)\pi/m\), \(0\le j<m\). On \([-1,1]\) the
unique monic polynomial of degree \(m\) of least uniform norm is
\(2^{1-m}T_m\), whose roots are \(\cos\theta_j\).
Interpolation at these \(m=n+1\) roots has
\(\Lambda_n=O(1+\log(n+1))\), with an absolute constant.
The affine images of the nodes give the corresponding minimax
monic nodal polynomial and the same Lebesgue bound on any fixed
nondegenerate compact interval. On that interval these interpolants
converge uniformly for every Hölder continuous function.
\end{theorem}

\begin{proof}
The cosine addition formula gives \(T_0=1\), \(T_1=X\) and
\(T_{j+1}=2XT_j-T_{j-1}\). In particular \(T_m\) has
degree \(m\) and leading coefficient \(2^{m-1}\). Its uniform
norm is \(1\) and its \(m\) distinct roots are as stated.
The fixed Mathlib revision has these facts and extremal statements
in the \texttt{Polynomial.Chebyshev} namespace \cite{mathlib}.

For minimax uniqueness, use the \(m+1\) decreasing points
\(z_i=\cos(i\pi/m)\), \(0\le i\le m\), and write
\(v_i=\prod_{j\ne i}(z_i-z_j)^{-1}\). Their signs are \((-1)^i\).
The leading coefficient of a degree-\(\le m\) polynomial \(q\)
is \(\sum_i v_i q(z_i)\) by Lagrange interpolation.
Since \(T_m(z_i)=(-1)^i\), the sum of the \(|v_i|\) is
\(2^{m-1}\). For a monic \(q\),
\[
  1\le \sum_i |v_i|\lVert q\rVert_\infty
    =2^{m-1}\lVert q\rVert_\infty .
\]
Equality forces every \(q(z_i)=2^{1-m}(-1)^i\), since each
term must attain its signed upper bound. Uniqueness of interpolation
then gives \(q=2^{1-m}T_m\).

For the Lebesgue estimate put \(x=\cos\theta\) with
\(0\le\theta\le\pi\). At a non-node, the \(j\)-th cardinal
polynomial satisfies
\[
 |\ell_j(\cos\theta)|
  =\frac{|\cos(m\theta)|\sin\theta_j}
         {m|\cos\theta-\cos\theta_j|}
  \le \frac{\pi\,|\cos(m\theta)|}{m|\theta-\theta_j|}.
\]
Indeed \(|T_m'(\cos\theta_j)|=m/\sin\theta_j\).
Factor the cosine difference as twice the product of the
sines of the half-sum and half-difference, in absolute value.
For fixed \(\theta_j\) the sine of the half-sum on the range
\(0\le\theta\le\pi\) is at least
\(\min(\sin(\theta_j/2),\cos(\theta_j/2))\),
so \(\sin\theta_j\le2\sin((\theta+\theta_j)/2)\).
Also \(|\sin((\theta-\theta_j)/2)|\ge|\theta-\theta_j|/\pi\).
These inequalities give the displayed bound.

Choose the closest angular node \(\theta_{j_*}\), with distance
at most \(\pi/(2m)\). The mean value theorem gives
\(|\cos(m\theta)|\le m|\theta-\theta_{j_*}|\), so the closest
term contributes at most \(\pi\). For each \(k\ge1\) there
are at most two indices with \(|j-j_*|=k\), and their angular
distance is at least \((k-\tfrac12)\pi/m\). Hence at every non-node
the sum is at most
\[
 \pi+2\sum_{k=1}^{m-1}\frac{1}{k-\tfrac12}
      =O(1+\log m).
\]
At a node the sum is \(1\). The affine change of variable
preserves all cardinal values and thus the interpolation norm.
It also gives the minimax monic polynomial on \([a,b]\), since
for every monic \(q\) of degree \(m\) on that interval the map
\[
 t\longmapsto
 \left(\frac{2}{b-a}\right)^m q\!\left(\frac{a+b}{2}
                              +\frac{b-a}{2}t\right)
\]
is monic and scales all uniform norms by the same factor.

For the convergence assertion, use that affine change to work on
\([0,1]\). If \(|f(t)-f(s)|\le L|t-s|^\alpha\), \(0<\alpha\le1\),
the degree-\(\le n\) Bernstein polynomial is
\(B_nf(t)=\mathbb E f(Z/n)\) with \(Z\) binomial of parameters
\(n,t\). Since \(\mathbb E Z=nt\) and
\(\mathbb E(Z/n-t)^2=t(1-t)/n\), concavity of \(u^{\alpha/2}\)
gives
\[
 |B_nf(t)-f(t)|
 \le L\,\mathbb E|Z/n-t|^\alpha
 \le L(4n)^{-\alpha/2}\qquad (n\ge1).
\]
The binomial moment identities also appear in
\cite[\texttt{K03.16.bernstein}]{d03}. Inserting \(B_nf\) into
Theorem~\ref{p206:W52} gives a bound of order
\((1+\log(n+1))n^{-\alpha/2}\), which tends to zero.
\end{proof}

\begin{theorem}[Failure of equispaced interpolation for continuous functions]\label{p206:W54}
For \(n+1\) equally spaced nodes including both endpoints of a
fixed nondegenerate interval, and \(n\ge2\),
\[
  \Lambda_n\ge\frac{2^n}{2n(n+1)^2}.
\]
In particular there exists a single continuous real function
on the interval whose equispaced interpolation polynomials
have unbounded uniform norm.
\end{theorem}

\begin{proof}
Rescale the nodes for a fixed \(n\) to \(0,1,\ldots,n\) and
evaluate at \(x=1/2\). The product over all nodes satisfies
\[
 \prod_{k=0}^n|x-k|
   =\frac{n!}{2}\prod_{k=1}^n\left(1-\frac{1}{2k}\right)
   \ge\frac{n!}{2(n+1)},
\]
since \(1-1/(2k)\ge k/(k+1)\). Choose a central index \(j\)
maximizing \(\binom nj\). The maximum is at least
\(2^n/(n+1)\), by summing all binomial coefficients.
Also \(|j-1/2|\le n\) and
\(\prod_{k\ne j}|j-k|=j!(n-j)!\). Thus the single term
\(|\ell_j(1/2)|\) is at least the stated lower bound.

Return to the fixed interval. Each \(I_n\) is a bounded linear
operator on \(C[a,b]\), a Banach space, and its norm is
\(\Lambda_n\) by Theorem~\ref{p206:W52}. If every continuous
\(f\) had \(\sup_n\lVert I_nf\rVert_\infty<\infty\), uniform
boundedness would give \(\sup_n\lVert I_n\rVert<\infty\),
contrary to the lower bound. Here the fixed Mathlib theorem is
\texttt{banach\_steinhaus} \cite{mathlib}. In its Baire argument
the closed sets \(\{f:\sup_n\lVert I_nf\rVert_\infty\le k\}\)
would cover the Banach space. A ball of radius \(r>0\) in one
of them then gives \(\lVert I_n g\rVert_\infty\le2k\) for
every \(\lVert g\rVert_\infty<r\), by subtracting the value
at its center. This yields a uniform operator norm bound.
\end{proof}
\chapter{Numerical Differentiation and Quadrature}
\label{p206:ch07}

Throughout this chapter the error of a quadrature rule \(Q\) is
the integral minus \(Q\). For \(f\) on a compact interval,
\(\lVert f\rVert_\infty\) denotes its uniform norm there. Numerical
differentiation will use separate estimates for the analytic
stencil and the errors in the sample values.

\section{Peano kernels and composite rules}

\begin{theorem}[Peano kernel representation]\label{p206:W55}
Let \(a<b\), let \(r\ge2\), and let \(E:C[a,b]\to\mathbb R\)
be a bounded linear functional vanishing on polynomials of
degree less than \(r\). Define \(u_+=\max(u,0)\) and
\[
 K(t)=E\!\left(\frac{(\,\cdot-t)_+^{r-1}}{(r-1)!}\right),
 \qquad a\le t\le b.
\]
Then \(K\) is continuous and every \(f\in C^r[a,b]\) satisfies
\[
 E(f)=\int_a^b f^{(r)}(t)K(t)\,dt,\qquad
 |E(f)|\le\lVert f^{(r)}\rVert_\infty\int_a^b|K(t)|\,dt.
\]
If \(K\) has a constant sign, its \(L^1\) norm is
\(\left|E((\,\cdot\,)^r/r!)\right|\).
\end{theorem}

\begin{proof}
Taylor's formula with integral remainder about \(a\) gives
\[
 f(x)=\sum_{k=0}^{r-1}\frac{f^{(k)}(a)}{k!}(x-a)^k
       +\int_a^b f^{(r)}(t)\frac{(x-t)_+^{r-1}}{(r-1)!}\,dt.
\]
The integral formula is proved in \cite[\texttt{K03.06.integral}]{d03}.
The kernel in the integral is jointly continuous on the compact
square since \(r\ge2\). Consequently \(t\mapsto
f^{(r)}(t)(\,\cdot-t)_+^{r-1}/(r-1)!\) is uniformly continuous
into \(C[a,b]\). The Riemann sums converge in the uniform norm
because for a mesh of size at most \(\delta\), the uniform error is
at most \(b-a\) times the modulus of continuity at \(\delta\).
Boundedness of \(E\) allows application to these sums and
passage to the limit. Its vanishing on the Taylor polynomial
leaves exactly the representation. The same uniform continuity
argument shows continuity of \(K\), and integration of absolute
values gives the error bound. Finally \(f(x)=x^r/r!\) has
\(f^{(r)}=1\), so \(E(f)=\int K\). Constant sign gives the
last assertion.
\end{proof}

For example, the difference between integration and a finite
weighted sum of point evaluations is bounded on \(C[a,b]\).
The representation in this setting is the Peano error formula
of \cite[Section 3.2, Theorem 3.2.3]{stoer}.

\begin{theorem}[Composite midpoint and trapezoidal errors]\label{p206:W56}
Let \(f\in C^2[a,b]\) with \(a<b\), \(n\ge1\), and
\(h=(b-a)/n\). Write \(M_n\) and \(T_n\) for the rules obtained
by adding the midpoint and trapezoidal rules on the \(n\)
subintervals of length \(h\). Then
\[
 \left|\int_a^b f-M_n(f)\right|
     \le \frac{(b-a)h^2}{24}\lVert f''\rVert_\infty,
 \qquad
 \left|\int_a^b f-T_n(f)\right|
     \le \frac{(b-a)h^2}{12}\lVert f''\rVert_\infty .
\]
\end{theorem}

\begin{proof}
Translate a subinterval to \([0,h]\). Both rules are exact
for constants and linear functions, by their zeroth and first
moments, so Theorem~\ref{p206:W55} applies with \(r=2\).
For \(0\le t\le h\) the integral of \((x-t)_+\) is
\((h-t)^2/2\). The midpoint rule subtracts \(h(h/2-t)_+\),
and its kernel becomes
\[
 K_M(t)=
 \begin{cases}
 t^2/2 & 0\le t\le h/2,\\
 (h-t)^2/2 & h/2\le t\le h.
 \end{cases}
\]
Its integral is \(2\int_0^{h/2}t^2/2\,dt=h^3/24\).
The trapezoidal rule subtracts \(h(h-t)/2\), leaving
\(K_T(t)=-t(h-t)/2\), whose integral is \(-h^3/12\).
Thus each local absolute error is bounded by the appropriate
kernel norm times the uniform second derivative bound.
The sum of the bounds uses \(nh^3=(b-a)h^2\).
\end{proof}

\begin{theorem}[Composite Simpson error]\label{p206:W57}
Let \(f\in C^4[a,b]\) with \(a<b\), let \(n\ge2\) be even, and
let \(h=(b-a)/n\). Form \(S_n\) by applying the three-point
weights \(h(1,4,1)/3\) to successive disjoint panels of length
\(2h\). Then
\[
 \left|\int_a^b f-S_n(f)\right|
      \le \frac{(b-a)h^4}{180}\lVert f^{(4)}\rVert_\infty .
\]
\end{theorem}

\begin{proof}
On \([0,2h]\), direct substitution into the rule for
\(1,x,x^2,x^3\) gives \(2h,2h^2,8h^3/3,4h^4\),
respectively, their exact integrals. Therefore the Peano
representation uses \(r=4\). On \(h\le t\le2h\), only the
right node contributes to the rule on \((x-t)_+^3/6\), so
\[
 K(t)=\frac{(2h-t)^4}{24}-\frac{h(2h-t)^3}{18}
      =\frac{(2h-t)^3(2h-3t)}{72}\le0 .
\]
Reflection \(x\mapsto2h-x\) preserves \(E\) and takes
\((x-(2h-t))_+^3\) to \((t-x)_+^3\).
The difference of \((x-t)_+^3\) and \((t-x)_+^3\) is
the cubic \((x-t)^3\), which \(E\) annihilates.
Hence \(K(2h-t)=K(t)\), so the kernel is nonpositive on
the whole panel. Its mass equals the error on \(x^4/24\).
\[
 \int_0^{2h}K(t)\,dt
 =\left(\frac4{15}-\frac5{18}\right)h^5=-\frac{h^5}{90}.
\]
Sum the resulting local bounds over \(n/2\) panels.
The composite rules and their error formulas are discussed in
\cite[Sections 3.1 and 3.2]{stoer}.
\end{proof}

\section{Gaussian quadrature and the Jacobi matrix}

Let \(\mu\) be a finite positive Borel measure on \(\mathbb R\)
concentrated on an interval. A finite moment of order \(2n\)
means \(\int |x|^{2n}\,d\mu<\infty\). All moments of lower
orders then exist absolutely. We use its closed support,
consisting of points all of whose open neighborhoods have
positive measure. Its complement is null because that open
complement is a countable union of open intervals of measure
zero, by the countable rational interval basis.

\begin{theorem}[Gaussian quadrature for a positive measure]\label{p206:W58}
Let \(n\ge1\). Assume that \(\mu\) has a finite moment of
order \(2n\) and that the Gram matrix
\[
 G_n=\left(\int x^{i+j}\,d\mu(x)\right)_{0\le i,j\le n}
\]
is positive definite. Equivalently, the support of \(\mu\)
contains at least \(n+1\) points. There is a unique monic
polynomial \(\pi_n\) of degree \(n\) orthogonal to every
polynomial of degree less than \(n\), and it has \(n\)
distinct real roots \(x_1,\ldots,x_n\) strictly between the
infimum and supremum of the support. There are positive
weights \(w_i\) such that
\[
 \int p\,d\mu=\sum_{i=1}^n w_i p(x_i)
 \qquad\text{for every }p\in\mathcal P_{2n-1}(\mathbb R).
\]
No deterministic rule using at most \(n\) point values can be exact on all
polynomials of degree at most \(2n\).
\end{theorem}

\begin{proof}
For a vector \(c\), the form \(c^\mathsf{T}G_nc\) equals
the integral of the square of the associated polynomial.
If the support contains at least \(n+1\) points, a nonzero
polynomial of degree at most \(n\) is nonzero at some
support point. By continuity its square is bounded below
by a positive number on a neighborhood of positive measure.
The form is therefore positive for nonzero \(c\).
If the support has at most \(n\) points, the product of
the corresponding linear factors is a nonzero polynomial
of degree at most \(n\) vanishing almost everywhere.
This proves the equivalence, also for the zero measure.

Write \(\pi_n(X)=X^n+\sum_{j=0}^{n-1}a_jX^j\). The
\(n\) orthogonality equations against \(1,x,\ldots,x^{n-1}\)
have the invertible leading principal Gram block as their
coefficient matrix. This gives existence and uniqueness.

For the root count, let \(q\) be the product of the distinct
real roots of \(\pi_n\) of odd multiplicity lying strictly
between the support extremes. If there are fewer than \(n\),
then \(\deg q<n\), and \(\pi_n q\) has a constant sign on
the open interval between the extremes because its interior
roots all have even multiplicity. The sign extends weakly
to any finite ends. Moreover the distinct zeros of \(\pi_n q\)
are zeros of \(\pi_n\), so at most \(n\) real points.
At some support point the product is nonzero. Its absolute
value is bounded below on a neighborhood of positive measure,
and the constant sign on the support makes the integral
nonzero. Integrability follows from the moment assumption
because \(\deg(\pi_n q)\le2n-1\). This contradicts
orthogonality. There must therefore be \(n\) such roots,
and the degree shows they are all simple.

For the cardinal polynomials \(\ell_i\) at these \(n\)
nodes set \(w_i=\int\ell_i\,d\mu\). Interpolation shows
exactness through degree \(n-1\). If \(\deg p\le2n-1\),
divide \(p=\pi_n s+r\) with \(\deg r<n\) and
\(\deg s\le n-1\). The integral of \(\pi_n s\) is zero,
and its nodal values vanish, so exactness for \(r\) proves
the claim. Applying exactness to \(\ell_i^2\) gives
\(w_i=\int\ell_i^2\,d\mu>0\), by positivity of the Gram
form. Finally, run a proposed deterministic point-evaluation rule
on the zero polynomial. Its at most \(n\) queried points determine
a monic nodal product whose square vanishes at all of them but has
positive integral by the Gram condition. Run the same rule on that
square. Every answer on the zero-data path stays zero, so even if
the rule chooses points adaptively it takes the same path and gives
the same output. It cannot be exact for both inputs.
\end{proof}

The construction is the general orthogonal-polynomial Gaussian
rule of \cite[Section 4.1]{gautschi} and
\cite[Section 3.6]{stoer}. The hypothesis for nondegeneracy
through degree \(n\) gives positive squared norms even at degree
\(n\), giving the exactness limit in the theorem.

\begin{theorem}[Gaussian remainder and Jacobi computation]\label{p206:W59}
Under the hypotheses of Theorem~\ref{p206:W58}, suppose in addition
for the analytic remainder that the support is contained in
\([a,b]\) and \(f\in C^{2n}[a,b]\). Then, with
\(Q_n(f)=\sum_iw_i f(x_i)\),
\[
 \left|\int f\,d\mu-Q_n(f)\right|
 \le\frac{\lVert f^{(2n)}\rVert_\infty}{(2n)!}
          \int \pi_n^2\,d\mu .
\]
In fact the error is \(f^{(2n)}(\xi)\int\pi_n^2\,d\mu/(2n)!\)
for some \(\xi\in[a,b]\).

For the exact computation of nodes and weights, let \(p_k\),
\(0\le k\le n\), be orthonormal polynomials for \(\mu\) with
positive leading coefficients. The Jacobi matrix of order \(n\)
represents multiplication by \(x\) followed by orthogonal
projection onto the span of \(p_0,\ldots,p_{n-1}\). It is real
symmetric tridiagonal, with positive subdiagonal entries where
present. Its eigenvalues are the \(x_i\). If \(v_i\) is a unit
eigenvector for \(x_i\), with coordinates indexed from zero, then
\[
 w_i=\mu(\mathbb R)\,|(v_i)_0|^2.
\]
These matrix formulas use the moment assumptions alone.
\end{theorem}

\begin{proof}
For the analytic error, interpolate \(f,f'\) at each \(x_i\)
by a polynomial \(H\) of degree at most \(2n-1\) using
Theorem~\ref{p206:W51}. Since \(H\) and \(f\) have the same
nodal values, \(Q_n(f)=Q_n(H)=\int H\,d\mu\).
Put \(m_*=\min_{[a,b]} f^{(2n)}\) and
\(M_*=\max_{[a,b]} f^{(2n)}\).
The Hermite remainder gives the pointwise bounds
\[
  \frac{m_*}{(2n)!}\pi_n(x)^2
  \le f(x)-H(x)
  \le\frac{M_*}{(2n)!}\pi_n(x)^2.
\]
Integrate these derivative bounds against the positive measure.
This avoids choosing a measurable intermediate point. Since
\(\int\pi_n^2\,d\mu>0\), the scaled error is between
\(m_*\) and \(M_*\), and continuity gives the exact
intermediate derivative value. The absolute bound follows.

For the computation, Gram--Schmidt on \(1,X,\ldots,X^n\)
is valid by the Gram condition. Multiplication by \(X\)
is symmetric for the polynomial inner product, since
\(\langle Xp_j,p_i\rangle=\langle p_j,Xp_i\rangle\).
If \(i<j-1\), the second term is zero by orthogonality
because \(\deg(Xp_i)<j\). Symmetry also gives the zeros
for \(j<i-1\). Thus for \(0\le j<n\) there is the
three-term recurrence
\[
 Xp_j=\beta_{j+1}p_{j+1}+\alpha_jp_j+\beta_jp_{j-1},
 \quad
 \alpha_j=\langle Xp_j,p_j\rangle,\qquad
 \beta_{j+1}=\frac{\operatorname{lc}(p_j)}
                   {\operatorname{lc}(p_{j+1})}>0 ,
\]
with \(p_{-1}=0,\ \beta_0=0\) and \(\operatorname{lc}\)
the leading coefficient. The formula for \(\beta_{j+1}\)
comes from comparing leading coefficients. Symmetry
identifies the coefficient of \(p_{j-1}\) with \(\beta_j\).
The Jacobi matrix has diagonal \(\alpha_j\) and off-diagonal
entries \(\beta_j\) between rows \(j-1,j\) for \(1\le j<n\).
The recurrence at a root \(x_i\) of \(p_n\) shows that
\((p_j(x_i))_{0\le j<n}\) is its eigenvector for \(x_i\),
since the term in \(p_n\) disappears. The vector is nonzero
because \(p_0=\mu(\mathbb R)^{-1/2}\).

To normalize these vectors, the matrix \(U\) with entries
\(U_{ji}=\sqrt{w_i}\,p_j(x_i)\) satisfies
\[
 (UU^\mathsf{T})_{jk}=\sum_{i=1}^n w_i p_j(x_i)p_k(x_i)
                     =\int p_jp_k\,d\mu=\delta_{jk},
\]
by exactness through degree \(2n-2\). Being square, \(U\)
has unit orthogonal columns. The \(n\) roots are distinct,
so each eigenspace is one-dimensional, and any unit
eigenvector for \(x_i\) agrees with the corresponding
column up to sign. Its first coordinate is
\(\sqrt{w_i}\,\mu(\mathbb R)^{-1/2}\) up to that sign.
This gives the claimed weight formula, the exact
Jacobi-matrix interface in \cite[Section 4.1]{gautschi}.
\end{proof}

\section{Finite differences and extrapolation}

\begin{theorem}[Truncation and sample error]\label{p206:W60}
Fix \(x\in\mathbb R\) and \(h>0\). Suppose the sampled values
\(y(t)\) at the stencil points satisfy \(|y(t)-f(t)|\le\eta\).
If \(f\in C^3[x-h,x+h]\) and
\(M_3=\sup_{[x-h,x+h]}|f^{(3)}|\), then the central difference
satisfies
\[
 \left|\frac{y(x+h)-y(x-h)}{2h}-f'(x)\right|
       \le \frac{M_3h^2}{6}+\frac{\eta}{h}.
\]
If \(f\in C^2[x,x+h]\) and \(M_2=\sup_{[x,x+h]}|f''|\),
then the forward difference satisfies
\[
 \left|\frac{y(x+h)-y(x)}h-f'(x)\right|
       \le\frac{M_2h}{2}+\frac{2\eta}{h}.
\]
These bounds assume exact subtraction and division of the
sampled values.
\end{theorem}

\begin{proof}
For the central formula use Taylor through order two at \(x\)
for each of \(x+h,x-h\). The quadratic and constant terms
cancel. Each remaining Taylor error has magnitude at most
\(M_3h^3/6\), giving \(M_3h^2/6\) after subtraction and
division by \(2h\). The two sample errors contribute at most
\(2\eta/(2h)\). For the forward formula Taylor through
order one has error at most \(M_2h^2/2\). Divide by \(h\)
and add the two scaled sample errors to get the second bound.
See \cite[\texttt{K03.06.lagrange}]{d03} for the Lagrange
remainder and \cite[Section 5.5]{driscoll} for the
differentiation error setting.
\end{proof}

For step selection, take \(M_2,M_3,\eta\) independent of \(h\)
on \(0<h\le H\) by using derivative bounds on the entire
allowed neighborhood and a uniform sample error bound. If both
terms in the central bound have positive coefficients, the
unconstrained minimum occurs at \(h=(3\eta/M_3)^{1/3}\),
since the derivative of the bound is \(M_3h/3-\eta/h^2\)
and changes from negative to positive there. For the forward
bound the derivative is \(M_2/2-2\eta/h^2\), giving
\(h=(4\eta/M_2)^{1/2}\). In each case cap this value by \(H\).
With positive sample error and zero derivative bound, the bound
decreases up to \(H\). With zero sample error and positive
derivative bound it tends to zero as \(h\) decreases to zero.
With both zero it vanishes identically. If forming the stencil
itself involves rounded subtraction or division, the arithmetic
model of Chapter~\ref{p206:ch01} supplies an additional error
term beyond the sample-error budget.

\pagebreak
\begin{proposition}[One-step Richardson extrapolation]\label{p206:W61}
Suppose \(A(h)=A+c h^p+O(h^{p+q})\) as \(h\downarrow0\),
where \(p,q>0\), and fix \(\lambda>1\). Then
\[
 R(h)=\frac{\lambda^p A(h/\lambda)-A(h)}{\lambda^p-1}
       = A+O(h^{p+q}).
\]
If the two approximations at \(h/\lambda\) and \(h\) have
absolute data errors at most \(\delta_{h/\lambda}\) and
\(\delta_h\), the same linear combination in exact arithmetic
has an additional error at most
\[
 \frac{\lambda^p\delta_{h/\lambda}+\delta_h}{\lambda^p-1}.
\]
For the central derivative stencil on exact data and \(f\in C^5\),
the choice \(p=2,\ \lambda=2\) gives a fourth-order formula.
\end{proposition}

\begin{proof}
The identity \(\lambda^p(h/\lambda)^p=h^p\) cancels the leading
error term. The constants leave \(A\), and the remainder is
\(O(h^{p+q})\) for fixed \(\lambda\). The input error bound follows
by the triangle inequality. For the central stencil, Taylor
through order four gives
\[
 \frac{f(x+h)-f(x-h)}{2h}
       =f'(x)+\frac{f^{(3)}(x)}6 h^2+r(h),\qquad
 |r(h)|\le\frac{\lVert f^{(5)}\rVert_\infty}{120}h^4,
\]
The even terms cancel, and each Taylor remainder is at most
\(\lVert f^{(5)}\rVert_\infty h^5/120\) on the stencil interval.
Thus \((4A(h/2)-A(h))/3\) has truncation error bounded by
\(\lVert f^{(5)}\rVert_\infty h^4/288\).
See \cite[Section 3.4]{stoer} for extrapolated quadrature.
\end{proof}
\chapter[Brackets, Fixed Points, and Stopping]{Bracketing, Fixed Points, and Certified Stopping}

A small residual and a small error are different quantities. A bracket
controls the location of a scalar root geometrically. A contraction
provides a different certificate, in which the residual is divided by a
positive distance from the limiting contraction factor \(1\). We first
develop both certificates in exact arithmetic, then specify what must be
verified when an evaluation is approximate. The scalar bracketing method
is discussed in \cite[Section~2.1]{persson}, and the fixed-point method in
\cite[Section~5.2]{stoer} and \cite[Section~4.2]{driscoll}.

\section{Brackets and safeguarded updates}

\begin{theorem}[Root existence in a bracket]\label{p206:W62}
Let \(a<b\) and let \(f\colon[a,b]\to\mathbb R\) be continuous. If
\(f(a)f(b)\le 0\), there is a root in \([a,b]\). If both endpoint values
are nonzero and have opposite signs, there is a root in \((a,b)\).
\end{theorem}
\begin{proof}
A zero endpoint is already a root. In the remaining case, zero lies
strictly between the endpoint values, so the intermediate value theorem
\cite[entry \texttt{K03.03.intermediate\_values}]{d03} gives a root
strictly between the endpoints. For example, \(x(x^2-1)\) has opposite
signs at \(-2\) and \(2\) and has all three roots \(-1,0,1\) in the
bracket.
\end{proof}

\begin{theorem}[Bisection and its bracket invariant]\label{p206:W63}
Under the assumptions of Theorem~\ref{p206:W62}, set
\(I_0=[a_0,b_0]=[a,b]\). For a current interval \(I_n\) form its exact
midpoint \(m_n=(a_n+b_n)/2\). Either return an endpoint or midpoint
known to be a root, or keep a half interval whose endpoint values
have product at most zero. For every step carried out,
\[
 I_{n+1}\subseteq I_n,\qquad f(a_n)f(b_n)\le 0,\qquad
 b_n-a_n=2^{-n}(b-a).
\]
Every midpoint is within \(2^{-n-1}(b-a)\) of a root in its bracket.
If the process continues indefinitely, the endpoints and midpoints
converge to a root \(x_*\), and the same bound holds for
\(\lvert m_n-x_*\rvert\) at every \(n\).
\end{theorem}
\begin{proof}
For strictly opposite endpoint signs, a nonzero midpoint value shares
the sign of one endpoint. Replacing that endpoint by the midpoint
keeps the product negative. A midpoint root can be returned or retained
as an endpoint of either half. If an endpoint already is a root, the
half containing it keeps the product at most zero. This proves the
invariant, and each completed update halves the interval width.
Theorem~\ref{p206:W62} supplies a root in every bracket, so its distance
from the midpoint is at most half the width.

For an infinite run the left endpoints are nondecreasing and the
right endpoints are nonincreasing. They are bounded, and their
difference tends to zero, so they have the same limit \(x_*\).
Continuity gives
\(f(x_*)^2=\lim_{n\to\infty}f(a_n)f(b_n)\le0\).
The limit is in every interval, giving the stated midpoint bound
relative to this one root.
\end{proof}

\begin{proposition}[Step count and certified signs]\label{p206:W64}
For \(\epsilon>0\), exact bisection requires at most
\[
 N_\epsilon=\max\left(0,\left\lceil
      \log_2\frac{b-a}{2\epsilon}\right\rceil\right)
\]
half-interval updates before its current midpoint certifies an absolute
root error of at most \(\epsilon\). Evaluating that final midpoint is
unnecessary if the bracket and the midpoint alone are the desired
certificate.

If a function evaluation at a point \(x\) returns a certified enclosure
\([l,u]\ni f(x)\), then \(l>0\) certifies a positive sign, \(u<0\)
certifies a negative sign, and \([l,u]=[0,0]\) certifies a root. A
strict-sign update requires a certified sign.
\end{proposition}
\begin{proof}
By Theorem~\ref{p206:W63}, half the width after \(n\) updates is
\((b-a)2^{-n-1}\). This is at most \(\epsilon\) exactly when
\(n\ge\log_2((b-a)/(2\epsilon))\), with \(n\) an integer and \(n\ge0\).
In particular, the initial midpoint suffices when the initial half-width
is at most \(\epsilon\). Each sign implication follows directly from
membership in \([l,u]\). If the enclosure contains zero without being
\([0,0]\), it certifies neither a strict sign nor a root. In that case
the existing bracket remains available, and a sharper enclosure may
allow another update. An interior representable point must also be
available when implementing a step with a fixed number format.
\end{proof}

With directed rounding, the enclosure hypotheses can be proved using
the containment properties from Chapter~1. The step-count formula
applies to completed exact halvings. In an implementation, if a stored
output \(z\) is inside a verified bracket \([a',b']\), any certified upper
bound on \(\max(z-a',b'-z)\) is a valid error radius for \(z\). If it
is impossible to represent the exact midpoint or to certify the sign at
a new point, the existing enclosure is still the certificate.

\begin{theorem}[Safeguarded bracketing]\label{p206:W65}
Start with a strict sign bracket and fix \(0<\alpha\le\tfrac12\).
If the current bracket is \([a_n,b_n]\) of width \(w_n\), a proposed
point from another method may be used when it is in the central
interval
\[
 [a_n+\alpha w_n,\ b_n-\alpha w_n].
\]
Otherwise use the exact midpoint. If the chosen point has a certified
nonzero sign, replace the endpoint with the same sign. If it is a
certified root, return it. Each successful update preserves a sign
bracket and satisfies \(w_{n+1}\le(1-\alpha)w_n\). Consequently an
infinite sequence of successful updates converges to a root, and
its \(n\)-th bracket midpoint has error at most
\((b-a)(1-\alpha)^n/2\) with respect to that root.
\end{theorem}
\begin{proof}
The accepted point is at distance at least \(\alpha w_n\) from
each endpoint. The midpoint satisfies the same condition.
A nonzero sign selects the opposite-sign subinterval, whose width
is at most \(w_n-\alpha w_n\). The nested interval argument of
Theorem~\ref{p206:W63} applies to this geometrically decreasing
width. If a proposed Newton or secant step is undefined, the midpoint
is an admissible choice. If the sign cannot be certified, or if
finite representation prevents selecting a point certified to be in
the central interval, the old bracket is retained. The stated
geometric bound counts only successful updates.
\end{proof}

\section{Exact and inexact fixed-point iterations}

\begin{theorem}[Banach's fixed-point theorem]\label{p206:W66}
Let \((X,d)\) be a nonempty complete metric space and let
\(T\colon X\to X\) satisfy
\[
 d(Tx,Ty)\le q\,d(x,y)\quad\text{for all }x,y\in X,\qquad
 0\le q<1.
\]
Then \(T\) has a unique fixed point \(p\), and \(x_n=T^n x_0\)
converges to \(p\) from every \(x_0\in X\).
\end{theorem}
\begin{proof}
Inductively, \(d(x_{j+1},x_j)\le q^j d(x_1,x_0)\), with \(q^0=1\)
also when \(q=0\). For \(m>n\), the triangle inequality yields
\[
 d(x_m,x_n)\le
 d(x_1,x_0)\sum_{j=n}^{m-1}q^j
 \le\frac{q^n}{1-q}d(x_1,x_0).
\]
The right-hand side tends to zero, so the iterates converge to
some \(p\in X\). The contraction inequality implies continuity of
\(T\). Passing to the limit in \(x_{n+1}=Tx_n\) gives \(Tp=p\).
For any other fixed point \(p'\), the inequality
\(d(p,p')\le qd(p,p')\) forces \(p=p'\).
This is the contraction argument also used in \cite{d10}.
\end{proof}

\begin{corollary}[A closed invariant domain]\label{p206:W67}
Let \(D\) be a nonempty closed subset of a complete metric space.
If \(T(D)\subseteq D\) and the contraction inequality holds for
points in \(D\), then there is a unique fixed point in \(D\), and the
iterates from any point of \(D\) converge to it.
\end{corollary}
\begin{proof}
With the restricted metric, every Cauchy sequence in \(D\) is Cauchy
in the ambient space, and its ambient limit lies in \(D\) by closedness.
Invariance makes the restriction of \(T\) a self-map of this complete
space, so Theorem~\ref{p206:W66} applies there. The iterates of this
restriction agree with the ambient iterates because all their points
stay in \(D\).
\end{proof}

\begin{proposition}[Sensitivity of a fixed point]\label{p206:W68}
Suppose \(T\) has fixed point \(p\), satisfies the contraction
inequality with \(0\le q<1\), and \(\widetilde T\) has a fixed point
\(\tilde p\) in the same space. If
\(d(Tx,\widetilde T x)\le\delta\) for all \(x\) in that space, then
\[
 d(p,\tilde p)\le\frac{\delta}{1-q}.
\]
In fact the perturbation bound at \(\tilde p\) suffices.
\end{proposition}
\begin{proof}
Use \(T\tilde p\) as the intermediate point in the triangle
inequality.
\[
 d(p,\tilde p)
 \le d(Tp,T\tilde p)+d(T\tilde p,\widetilde T\tilde p)
 \le qd(p,\tilde p)+\delta.
\]
Rearrangement gives the claim. To supply the existence premise,
one can prove that \(\widetilde T\) is itself a contraction on a
nonempty closed invariant subset of the complete domain. The resulting
fixed point then comes from Corollary~\ref{p206:W67}.
\end{proof}

\begin{proposition}[A priori and a posteriori bounds]\label{p206:W69}
For the iteration of Theorem~\ref{p206:W66}, the bounds
\[
 d(x_n,p)\le\frac{q^n}{1-q}d(x_1,x_0),\qquad
 d(x_n,p)\le\frac{d(x_{n+1},x_n)}{1-q}
\]
hold for all \(n\ge0\).
\end{proposition}
\begin{proof}
Take \(m\to\infty\) in the geometric-tail bound in the proof of
Theorem~\ref{p206:W66}. This gives the first inequality.
Starting that estimate with the new initial point \(x_n\) gives the
second. When \(q=0\), \(T\) is constant, hence \(x_1=p\) and the bounds
hold with \(q^0=1\). When \(x_0=p\), both right-hand sides are zero.
\end{proof}

\begin{proposition}[Fixed-point residual as a stopping certificate]\label{p206:W70}
For any \(x\) in a domain on which \(T\) is a \(q\)-contraction
with fixed point \(p\),
\[
 d(x,p)\le \frac{d(x,Tx)}{1-q}.
\]
In particular, a verified upper bound on \(d(x,Tx)\) of at most
\((1-q)\epsilon\) certifies \(d(x,p)\le\epsilon\).
\end{proposition}
\begin{proof}
The triangle inequality gives
\(d(x,p)\le d(x,Tx)+d(Tx,Tp)\le d(x,Tx)+q\,d(x,p)\).
Since \(1-q>0\), we can move the last term to the left and divide.
The conclusion continues to hold with any upper bound substituted
for the exact residual.
\end{proof}

\begin{theorem}[Inexact contraction iteration]\label{p206:W71}
In the setting of Theorem~\ref{p206:W66} or
Corollary~\ref{p206:W67}, let \(p\) be the fixed point.
Suppose an approximate sequence \(x_j\) remains in the contraction
domain and satisfies \(d(x_{j+1},Tx_j)\le\delta_j\) for all \(j\ge0\).
Then
\[
 d(x_n,p)\le q^n d(x_0,p)+
             \sum_{j=0}^{n-1}q^{n-1-j}\delta_j .
\]
For a common bound \(\delta_j\le\delta\) the upper limit of these
distances is at most \(\delta/(1-q)\). If \(\delta_j\to0\), then
\(x_n\to p\).
\end{theorem}
\begin{proof}
The triangle and contraction inequalities yield
\(d(x_{j+1},p)\le\delta_j+q\,d(x_j,p)\). Induction gives the
displayed finite sum, which is empty at \(n=0\). For a common
bound it is at most \(\delta(1-q^n)/(1-q)\). If \(\delta_j\to0\),
given \(\varepsilon>0\) choose \(J\) so that \(\delta_j\le
(1-q)\varepsilon/2\) for all \(j\ge J\). The corresponding tail
of the sum is at most \(\varepsilon/2\). The contribution of the
finitely many terms \(j<J\) and of the initial distance tends to
zero. This proves convergence, including \(q=0\).

If \(y\) is an approximation to \(Tx\) with certified error at most
\(\delta_x\), the triangle inequality gives
\(d(x,Tx)\le d(x,y)+\delta_x\).
Thus an upper bound on \(d(x,y)\) plus \(\delta_x\) can be used
in Proposition~\ref{p206:W70}. Verify separately that the computed
iterates stay in the contraction domain.
\end{proof}

\begin{proposition}[Inverse-slope certificate for a scalar root]\label{p206:W72}
Let \(x,x_*\) lie in an interval, with \(f(x_*)=0\). Assume \(f\)
is continuous on the closed segment between \(x\) and \(x_*\),
differentiable in its interior, and \(\lvert f'(t)\rvert\ge m>0\)
there. Then
\[
 \lvert x-x_*\rvert\le\frac{\lvert f(x)\rvert}{m}.
\]
\end{proposition}
\begin{proof}
When \(x=x_*\) the claim is immediate. Otherwise the real mean value
theorem \cite[entry \texttt{K03.05.lagrange}]{d03} gives
\(f(x)-f(x_*)=f'(\xi)(x-x_*)\) for an interior point \(\xi\).
Taking absolute values proves the estimate. The assumptions include
the usual continuously differentiable case. A verified lower slope bound
on an entire certified bracket containing both \(x\) and a root also
suffices, even when the root location is unknown.

For \(f(x)=x^2\) at its root zero the residual is the square of the
error, so no fixed linear inverse-residual bound holds on a
neighborhood of that root. Even at simple roots, a bound uniform
in the function cannot follow from a small residual alone.
The functions \(f_\sigma(x)=\sigma x\) with \(\sigma>0\) have
residual \(\sigma\) at \(x=1\) and root error \(1\).
\end{proof}
\chapter{Local and Semilocal Newton Methods}

Local theorems assume that a root is already known to exist and
describe a neighborhood from which Newton's method converges.
A semilocal theorem uses data at the proposed starting point and
bounds on a specified ball to prove existence as well as convergence.
This distinction is important when the goal is a certificate.

\section{Scalar Newton and secant iterations}

\begin{theorem}[Local quadratic convergence at a simple root]\label{p206:W73}
Let \(f\) be real and \(C^2\) on a neighborhood of
\([x_*-R,x_*+R]\), where \(R>0\) and \(f(x_*)=0\). Suppose that on this
interval \(\lvert f'\rvert\ge m>0\) and \(\lvert f''\rvert\le M\).
Let \(0<r\le R\) with \(Mr/(2m)<1\). From
\(\lvert x_0-x_*\rvert\le r\), the exact Newton iterates
\(x_{n+1}=x_n-f(x_n)/f'(x_n)\) stay in the closed radius-\(r\) interval
and converge to \(x_*\). With \(e_n=x_n-x_*\),
\[
 \lvert e_{n+1}\rvert\le\frac{M}{2m}\lvert e_n\rvert^2 .
\]
If \(M=0\) then one step suffices.
\end{theorem}
\begin{proof}
Apply Taylor's formula at \(x_n\) with endpoint \(x_*\),
\[
 0=f(x_n)+f'(x_n)(x_*-x_n)+R_n,\qquad
 \lvert R_n\rvert\le\tfrac12 M\lvert e_n\rvert^2 .
\]
The formula is the real Lagrange remainder of \cite{d03},
entry \texttt{K03.06.lagrange}. Since \(f'(x_n)\ne0\),
rearrangement gives \(e_{n+1}=R_n/f'(x_n)\). In the radius-\(r\)
interval this implies
\(\lvert e_{n+1}\rvert\le (Mr/(2m))\lvert e_n\rvert\).
Induction gives invariance and convergence. The remainder vanishes
when \(M=0\). A root that is reached remains fixed under the iteration.
See also \cite[Section~2.3]{persson}.
\end{proof}

\begin{proposition}[Newton iteration at a multiple root]\label{p206:W74}
Near \(x_*\), suppose
\(f(x)=(x-x_*)^m g(x)\), where \(m\ge2\) is an integer,
\(g\in C^2\), and \(g(x_*)\ne0\). The usual Newton map, defined
away from the root in a sufficiently small neighborhood, satisfies
\[
 e_{\mathrm{new}}=(1-1/m)e+O(e^2),\qquad e=x-x_*.
\]
The corrected step \(x-mf(x)/f'(x)\) satisfies
\(e_{\mathrm{new}}=O(e^2)\). Both maps give convergent iterations
from a sufficiently small neighborhood, with an iterate at the
root interpreted as a successful stop. For a nonterminating
ordinary Newton iteration the absolute error ratio tends to
\(1-1/m\).
\end{proposition}
\begin{proof}
For \(e\ne0\) the derivative is
\(e^{m-1}(m g(x)+e g'(x))\). The quantity in parentheses tends to
\(m g(x_*)\ne0\), so its absolute value has a positive lower
bound near \(x_*\). Cancelling the common factor gives the exact
identity
\[
 \frac{f(x)}{f'(x)}
 =\frac{e\,g(x)}{m g(x)+e g'(x)}
 =\frac{e}{m}-\frac{e^2 g'(x)}{m(m g(x)+e g'(x))}.
\]
The last coefficient is bounded near \(x_*\), which gives both
error formulas. For the ordinary map choose a radius so that the
absolute linear factor plus the quadratic remainder divided by
\(\lvert e\rvert\) is strictly below \(1\). For the corrected map
choose a radius so that its quadratic estimate also implies a
strict contraction of the distance to the root. Thus in each case
the errors stay within the chosen neighborhood and tend to zero.
The first formula then gives the error ratio for the ordinary map.
This correction uses the specified multiplicity
\cite[Section~2.4]{persson}.
\end{proof}

\begin{theorem}[Local convergence and order of the secant method]\label{p206:W75}
Let \(f\) be real and \(C^2\) near a simple root \(x_*\).
Take two distinct initial points \(x_0,x_1\) sufficiently close to
\(x_*\). Stop at any root encountered. Otherwise the exact
secant iterates
\[
 x_{n+1}=x_n-\frac{f(x_n)}{f[x_{n-1},x_n]}
 \qquad (n\ge1)
\]
are well defined and converge to \(x_*\). Here \(f[u,v]\) is
the first divided difference. If the iteration does not terminate
and \(f''(x_*)\ne0\), let
\[
 \varphi=\frac{1+\sqrt5}{2},\qquad
 C=\left|\frac{f''(x_*)}{2 f'(x_*)}\right|.
\]
Then
\(\lim_{n\to\infty}\lvert e_{n+1}\rvert/\lvert e_n\rvert^\varphi
 =C^{1/\varphi}>0\).
\end{theorem}
\begin{proof}
Choose an interval about \(x_*\) on which \(\lvert f'\rvert\ge m>0\)
and \(\lvert f''\rvert\le M\). The mean value theorem bounds each
nondegenerate first divided difference below in magnitude by \(m\).
For distinct \(x_{n-1},x_n,x_*\), consider the affine interpolant
\(\ell\) through \(x_{n-1},x_n\). Its root is \(x_{n+1}\), and
\(\ell(x_*)=-f[x_{n-1},x_n]e_{n+1}\).
The divided-difference interpolation identity
(Theorem~\ref{p206:W48}) gives
\[
 f(x_*)-\ell(x_*)
   = f[x_{n-1},x_n,x_*](x_*-x_{n-1})(x_*-x_n).
\]
Since \(f(x_*)=0\), we obtain
\[
 e_{n+1}
 =\frac{f[x_{n-1},x_n,x_*]}{f[x_{n-1},x_n]}\,e_n e_{n-1}.
\]
Equivalently, the affine interpolation remainder
(Theorem~\ref{p206:W50}) identifies the second divided difference
with \(f''(\xi)/2\) for some \(\xi\) in the convex hull of the three
points. In particular
\(\lvert e_{n+1}\rvert\le (M/(2m))\lvert e_n e_{n-1}\rvert\).
Choose a radius \(r\) within the interval with \(Mr/(2m)<1\).
If both errors have absolute value at most \(r\), this estimate
keeps the next point in the interval and decreases the error from
the current point by at least a fixed geometric factor. Before
termination, \(f(x_n)\ne0\), so \(x_{n+1}\ne x_n\); hence each
successive divided difference continues to be defined. If \(M=0\),
the same estimate gives an exact root at the first secant step.

For a nonterminating run write \(E_n=\lvert e_n\rvert\) and
\[
 C_n=\left|\frac{f[x_{n-1},x_n,x_*]}{f[x_{n-1},x_n]}\right|.
\]
The mean value representations and continuity show \(C_n\to C\).
Each \(C_n>0\), since \(E_{n+1}=C_nE_nE_{n-1}\) and all the errors
are positive. For \(r_n=E_n/E_{n-1}^{\varphi}\), the identity
\(\varphi(\varphi-1)=1\) gives
\[
 r_{n+1}=C_n r_n^{-1/\varphi}.
\]
Thus \(s_n=\log r_n\) satisfies
\(s_{n+1}=\log C_n-s_n/\varphi\). Subtract
\(s=(\log C)/\varphi\), the fixed point of the limiting recursion.
The resulting recurrence has coefficient of magnitude \(1/\varphi<1\)
and a forcing term tending to zero. Unfolding it gives a decaying
initial term and a geometric convolution of the forcing terms.
Splitting the convolution at a fixed index shows it tends to zero,
as in Theorem~\ref{p206:W71}. Hence \(s_n\to s\), proving the asserted
error order. Compare \cite[Section~2.3]{persson} and
\cite[Section~4.4]{driscoll}.
\end{proof}

\section{Banach-space Newton estimates}

We take the derivative to be the Fr\'echet derivative. For complex
spaces the line parameter in the following argument is real. A ball
on which estimates are imposed will be contained in an open domain
of differentiability.

\begin{lemma}[Remainder on a segment]\label{p206:newton-segment}
Let \(F\colon U\to Y\) be differentiable, with \(U\) open in a
Banach space \(X\) and \(Y\) Banach. Let \(D\subset U\) be convex and
suppose \(F'\) is Lipschitz with constant \(L\) on \(D\).
For \(x,y\in D\),
\[
 \|F(y)-F(x)-F'(x)(y-x)\|\le\tfrac12 L\|y-x\|^2 .
\]
\end{lemma}
\begin{proof}
For \(d=y-x\), the segment path \(s\mapsto F(x+sd)\) has continuous
derivative \(F'(x+sd)d\) on \([0,1]\). Its integral is
\(F(y)-F(x)\). Indeed the continuous Banach-valued derivative has a
Riemann integral by completeness, and the scalar fundamental theorem
applied after each bounded linear functional gives this identity,
since bounded functionals commute with Riemann sums and separate
points by Hahn--Banach. Subtracting \(F'(x)d\) inside the integral
bounds the integrand in norm by \(Ls\|d\|^2\), giving the
factor \(1/2\).
\end{proof}

\begin{theorem}[Local Newton in Banach spaces]\label{p206:W76}
Let \(F\colon U\to Y\) be differentiable on an open subset of
a Banach space \(X\), with \(Y\) Banach. Suppose \(F(x_*)=0\),
\(F'(x_*)\colon X\to Y\) is a bounded isomorphism, and \(F'\) is
Lipschitz with constant \(L\) on \(\overline B(x_*,R)\subset U\).
Set \(b=\|F'(x_*)^{-1}\|\). For \(0<r\le R\) with
\[
 bLr<1,\qquad \frac{bLr}{2(1-bLr)}<1 ,
\]
Newton's method from any \(x_0\in\overline B(x_*,r)\) is defined,
stays in this ball, and converges to \(x_*\) with the bound
\[
 \|x_{n+1}-x_*\|\le \frac{bL}{2(1-bLr)}\|x_n-x_*\|^2.
\]
\end{theorem}
\begin{proof}
Such an \(r\) exists by taking it sufficiently small.
For \(x\) in that ball,
\(\|F'(x_*)^{-1}(F'(x)-F'(x_*))\|\le bLr<1\).
The Neumann-series inverse estimate (Theorem~\ref{p206:W10})
gives invertibility of \(F'(x)\) and
\(\|F'(x)^{-1}\|\le b/(1-bLr)\). At the Newton point
\(N_F(x)=x-F'(x)^{-1}F(x)\), use \(F(x_*)=0\) to write
\[
 N_F(x)-x_*
  =F'(x)^{-1}\bigl[F(x_*)-F(x)-F'(x)(x_*-x)\bigr].
\]
Lemma~\ref{p206:newton-segment} proves the error estimate. With the
chosen \(r\), this bounds the next error by a fixed fraction of the
current error, proving invariance and convergence. When \(L=0\),
the bound gives one-step convergence. A finite-dimensional account
appears in \cite[Section~5.3]{stoer}.
\end{proof}

\section{The Newton--Kantorovich theorem}

\begin{theorem}[Semilocal radius, error, and uniqueness]\label{p206:W77}
Let \(X,Y\) be Banach spaces, let \(U\subseteq X\) be open, and
let \(F\colon U\to Y\) be differentiable. At a starting point
\(x_0\in U\) assume that \(F'(x_0)\) is a bounded isomorphism,
and put \(B=F'(x_0)^{-1}\). Suppose \(\eta,L>0\), \(2\eta L\le1\),
and
\[
 \|BF(x_0)\|\le\eta,\qquad
 \|B(F'(x)-F'(y))\|\le L\|x-y\|
 \quad (x,y\in\overline B(x_0,t_*)),
\]
where \(\overline B(x_0,t_*)\subset U\) and
\[
 D=\sqrt{1-2\eta L},\qquad
 t_*=\frac{1-D}{L},\qquad t_{**}=\frac{1+D}{L}.
\]
For \(m(t)=\eta-t+\tfrac12 Lt^2\), define
\(t_0=0\) and \(t_{n+1}=t_n-m(t_n)/m'(t_n)\).
Then the Newton sequence \(x_{n+1}=N_F(x_n)\) is well defined,
\(\|x_n-x_0\|\le t_n<t_*\), and it converges to the unique zero
\(x_*\) of \(F\) in the closed ball \(\overline B(x_0,t_*)\).
The estimates
\[
 \|x_{n+1}-x_n\|\le t_{n+1}-t_n,\qquad
 \|x_*-x_n\|\le t_*-t_n,\qquad
 \|x_*-x_{n+1}\|\le \tfrac12\|x_*-x_n\|
\]
hold also when \(2\eta L=1\). If \(2\eta L<1\), the last estimate
can be supplemented by
\[
 \|x_*-x_{n+1}\|\le\frac{L}{2D}\|x_*-x_n\|^2 .
\]
If \(2\eta L<1\) and the domain and derivative bound extend
to \(\overline B(x_0,\rho)\) for \(t_*\le\rho<t_{**}\), the zero
is unique on that larger closed ball as well.
\end{theorem}
\begin{proof}
\emph{Scalar majorant.}
For \(0\le t<t_*\), we have \(m(t)>0\) and \(1-Lt=-m'(t)>0\).
Write \(s=m(t)/(1-Lt)>0\). Expanding the quadratic at \(t\) gives
\[
 \begin{split}
 m(t+s)&=\tfrac12 Ls^2,\\
 t_*-(t+s)&=\frac{L}{2(1-Lt)}(t_*-t)^2 .
 \end{split}
\]
The second identity uses \(m(t_*)=0\) and implies \(t<t+s<t_*\).
Consequently the scalar iterates increase and are bounded by \(t_*\).
If their limit were below \(t_*\), continuity of the scalar update
there would force \(m\) to vanish at that limit. Thus \(t_n\uparrow t_*\).
Writing \(d_n=t_*-t_n>0\) and
\(c_n=L/(2(1-Lt_n))\) gives \(d_{n+1}=c_n d_n^2\).
Since \(1-Lt_n=D+Ld_n\), we have
\(c_nd_n\le\tfrac12\), with equality at the critical case \(D=0\).
For \(D>0\), \(c_n\le L/(2D)\).

\emph{Majorization of the Newton steps.}
Set \(G=BF\colon U\to X\). Then \(G'(x_0)=I\) and the stipulated
Lipschitz constant for \(G'\) is \(L\).
Where \(G'\) is invertible, \(F'\) is invertible and
\(F'(x)^{-1}F(x)=G'(x)^{-1}G(x)\), so both Newton maps agree.
At step \(n\), assume
\(\|x_n-x_0\|\le t_n\) and \(\|G(x_n)\|\le m(t_n)\), as holds at
\(n=0\). The estimate
\(\|G'(x_n)-I\|\le Lt_n<1\) and Theorem~\ref{p206:W10} yield
\[
 \|G'(x_n)^{-1}\|\le\frac1{1-Lt_n},\qquad
 \|x_{n+1}-x_n\|\le\frac{m(t_n)}{1-Lt_n}=t_{n+1}-t_n .
\]
The triangle inequality gives
\(\|x_{n+1}-x_0\|\le t_{n+1}<t_*\). In particular the segment between
the two iterates lies in the required ball. The defining Newton
equation cancels the first-order terms in
Lemma~\ref{p206:newton-segment}, giving
\[
 \|G(x_{n+1})\|
 \le \tfrac12 L\|x_{n+1}-x_n\|^2
 \le \tfrac12 L(t_{n+1}-t_n)^2
 = m(t_{n+1}).
\]
This closes the induction. For \(j>n\), summing step bounds
gives \(\|x_j-x_n\|\le t_j-t_n\). The sequence is Cauchy, and
completeness gives a limit \(x_*\) in the closed ball with
\(\|x_*-x_n\|\le d_n\). Since \(m(t_n)\to0\), continuity of \(G\)
gives \(G(x_*)=0\) and hence \(F(x_*)=0\).

\emph{Rates and uniqueness.}
For any zero \(y\) in \(\overline B(x_0,t_*)\), the error identity
from Theorem~\ref{p206:W76}, applied to \(G\) on the segment from
\(x_n\) to \(y\), gives
\[
 \|y-x_{n+1}\|\le c_n\|y-x_n\|^2 .
\]
At \(n=0\) we have \(\|y-x_0\|\le d_0\).
Induction with \(d_{n+1}=c_nd_n^2\) shows \(\|y-x_n\|\le d_n\),
so \(y=x_*\). For \(y=x_*\) the same inequality and
\(c_nd_n\le1/2\) give the stated linear estimate, while the uniform
bound on \(c_n\) for \(D>0\) gives the quadratic estimate.
All inverses used here are taken at finite iterates with
\(Lt_n<1\), including when \(D=0\). For example,
\(F(t)=\eta-t+Lt^2/2\) with \(2\eta L=1\) has a double zero at
\(t_*\). Its Newton error halves at each step from zero.

Finally, if \(D>0\) and the assumptions hold on the larger ball,
take a zero \(y\) there and set \(s=\|y-x_0\|\le\rho\).
Taylor's estimate at \(x_0\) gives
\[
 s-\eta\le\|G(x_0)+(y-x_0)\|\le\tfrac12 Ls^2 .
\]
Thus \(m(s)\ge0\). The two distinct roots of \(m\) are \(t_*\) and
\(t_{**}\), and \(m<0\) strictly between them. Since \(s<t_{**}\),
we have \(s\le t_*\), reducing to the uniqueness already proved.
In the scalar quadratic example with \(2\eta L<1\), \(t_{**}\)
is a second zero, so the open upper bound on the radius is
necessary in general. The majorant argument is based on
\cite[Theorem~1 and Corollaries~10--12]{ferreira}.
\end{proof}

\begin{proposition}[Degenerate data and a stable radius expression]\label{p206:W78}
In the semilocal setting, if the residual bound is \(\eta=0\),
then \(x_0\) is already a zero and its certified existence
radius is zero. If \(L=0\), suppose the radius-\(\eta\) closed
ball lies in the domain and the derivative bound with \(L=0\)
holds there. Then a single Newton step gives a zero in that
ball, and the zero is unique there. For \(L>0\) and
\(0\le 2\eta L\le1\), the exact identity
\[
 t_*=\frac{2\eta}{1+\sqrt{1-2\eta L}}
\]
avoids subtracting two nearly equal numbers in the numerator.
\end{proposition}
\begin{proof}
If \(\eta=0\), then \(BF(x_0)=0\). Invertibility of \(B\) implies
\(F(x_0)=0\). If \(L=0\), the derivative of \(G=BF\) is \(I\)
throughout the specified ball. Applying
Lemma~\ref{p206:newton-segment} with zero constant, for any
\(d\) in that ball centered at the origin we obtain
\(G(x_0+d)=G(x_0)+d\). The Newton displacement from \(x_0\) is
\(-G(x_0)\), whose norm is at most \(\eta\). Its endpoint is the
unique zero by the affine formula. This also covers the
intersection of the two degenerate cases. For \(L>0\), multiply
the numerator \(1-\sqrt{1-2\eta L}\) by its conjugate and divide
by \(L(1+\sqrt{1-2\eta L})\). The resulting formula extends to
\(L=0\) with value \(\eta\). For a numerical certificate the
data and the computed radius still require certified bounds.
\end{proof}

\backmatter
\begingroup
\small
\begin{thebibliography}{99}
\setlength{\itemsep}{0.25em}
\bibitem{higham} N. J. Higham, \emph{Accuracy and Stability of
  Numerical Algorithms}, 2nd ed., SIAM, 2002.
\bibitem{golub} G. H. Golub and C. F. Van Loan, \emph{Matrix
  Computations}, 4th ed., Johns Hopkins University Press, 2013.
\bibitem{stoer} J. Stoer and R. Bulirsch, \emph{Introduction to
  Numerical Analysis}, 3rd ed., Springer, 2002.
\bibitem{driscoll} T. A. Driscoll and R. J. Braun, \emph{Fundamentals
  of Numerical Computation}, SIAM and the
  \href{https://fncbook.com/}{authors' online edition}.
\bibitem{trefpde} L. N. Trefethen, \emph{Finite Difference and
  Spectral Methods for Ordinary and Partial Differential Equations},
  \href{https://people.maths.ox.ac.uk/trefethen/4all.pdf}{Chapter 4, Section 4.2}.
\bibitem{berrut} J.-P. Berrut and L. N. Trefethen, Barycentric
  Lagrange interpolation, \emph{SIAM Review} 46 (2004), 501--517.
  \href{https://doi.org/10.1137/S0036144502417715}{DOI}.
\bibitem{trefapprox} L. N. Trefethen, \emph{Approximation Theory
  and Approximation Practice}, SIAM, 2013, Chapter 15,
  \href{https://doi.org/10.1137/1.9781611975949.ch15}{Lebesgue constants}.
\bibitem{gautschi} W. Gautschi,
  \href{https://www.cs.purdue.edu/homes/wxg/selected_works/section_11/179.pdf}
       {Orthogonal Polynomials, Quadrature, and Approximation:
        Computational Methods and Software (in Matlab)},
  author manuscript, Section 4.1, pp.~36--40.
\bibitem{persson} P.-O. Persson,
  \href{https://persson.berkeley.edu/math128a/lecture_notes/math128a_lecture_notes_chapter_2.html}
       {Solutions of Equations in One Variable},
  Math 128A lecture notes, Chapter 2.
\bibitem{ferreira} O. P. Ferreira and B. F. Svaiter,
  \emph{Kantorovich's Theorem on Newton's Method},
  \href{https://arxiv.org/abs/1209.5704}{arXiv:1209.5704}, 2012.
\bibitem{hoffman} A. J. Hoffman and H. W. Wielandt,
  The variation of the spectrum of a normal matrix,
  \emph{Duke Math. J.} 20 (1953), 37--39.
  \href{https://doi.org/10.1215/S0012-7094-53-02004-3}{DOI}.
\bibitem{mirsky} L. Mirsky, Symmetric gauge functions and unitarily
  invariant norms, \emph{Quart. J. Math.} 11 (1960), 50--59.
  \href{https://doi.org/10.1093/qmath/11.1.50}{DOI}.
\bibitem{d03} MathGenome, D03 unified notes on single-variable analysis,
  Modules 3, 5, 6, and 16.
\bibitem{d04} MathGenome, D04 unified notes on multivariable analysis,
  Chapter 2, segment mean value bound.
\bibitem{d05} MathGenome, \emph{Linear Algebra and Foundational Higher
  Algebra}, D05 unified notes, Modules 9, 13, and 14.
\bibitem{d07} MathGenome, D07 unified notes on abstract algebra,
  Module 9, finite Chinese remainder theorem.
\bibitem{d10} MathGenome, D10 point-set topology,
  Chapter 14, Banach contraction theorem.
\bibitem{p205} MathGenome, \emph{Extremal Methods, Networks, and
  Matroid Optimization}, unified notes P2-05, Chapter 8, W42.
\bibitem{mathlib} The Mathlib Community, \emph{Mathlib},
  \href{https://github.com/leanprover-community/mathlib4/tree/56a005855ea17e0b0958f73ec801ae140b0a1679}
       {fixed revision \texttt{56a005855ea1}}.
\end{thebibliography}
\endgroup
\end{document}
